% Tensor-network TikZ styles adapted from
% Papers/1804.04964/paper_normal.tex (Andras Molnar et al.)
% and kept visually close to the diagram language used there and in
% Papers/2011.12127/TN-Review-main.tex.
%
% Layout policy:
% - Public \TN... macros are chapter-facing and semantically named.
% - Private \TN@... constructions in tex/tn/tn_core.tex and
%   tex/tn/tn_library.tex provide the atomic grammar and reusable networks.
% - Public diagrams default to unlabeled tensor glyphs; the surrounding
%   formula/prose carries the tensor names.
% - A local tensor is a black dot.  A suppressed contracted component is a
%   neutral square box, a linear map is a rounded blue box, a state is a
%   rounded gray box, a formal expression has its own neutral box, and an
%   endomorphism inserted on an index is a red dot.
% - Solid strokes are contracted indices.  Dashed strokes indicate only a
%   grouping boundary, a change of multiplicity basis, or an identified
%   boundary of a periodic lattice.
% Repository-wide tensor-network calculus.  The input path is local to this
% entry point; clients load the sole public TikZ library by name.
\makeatletter
\def\input@path{{../../tex/tn/}}
\makeatother
\usetikzlibrary{tn}
 % chktex 27 (job-dir-relative path; standalone lint)

\makeatletter

% ---------- Public semantic macros ----------

\newcommand{\TNTikZDiagram}[2]{#2}

\newcommand{\TNMPSLocal}[2]{%
    \begin{tikzpicture}[tn picture]
        \TN@mpssiteWithOpenLegs{tnA}{(0,0)}{#1}
        \TN@label{tnANOpen}{(0,0.22)}{#2}
    \end{tikzpicture}%
}

\newcommand{\TNMPSWord}[4]{%
    \begin{tikzpicture}[tn picture]
        \TN@openMPSword{tn}{#2}{#3}{#4}
        \node[anchor=east] at (-0.85,0) {$#1$};
    \end{tikzpicture}%
}

\newcommand{\TNMPV}[4]{%
    \begin{tikzpicture}[tn picture]
        \TN@blockedMPSword{tn}{#2}{#3}
        \node[anchor=east] at (-0.85,0) {$#1$};
        \node at (\TN@chainellipsisx,-0.70) {$#4$};
    \end{tikzpicture}%
}

% Blocking: L neighbouring local tensors grouped (dashed box) into one
% blocked tensor; physical legs poke out the top, the outer virtual legs
% become the blocked bond.
\newcommand{\TNBlocking}[4]{%
    \begin{tikzpicture}[tn picture]
        \TN@openMPSword{tn}{#2}{#3}{#4}
        \TN@grouppath{
            ($(tnL.west)+(-0.30,0.22)$)
            rectangle ($(tnR.east)+(0.30,-0.34)$);}
    \end{tikzpicture}%
}

% MPV overlap: the top row carries copies of #1, the bottom row the conjugate
% copies of #2; the shared physical indices are contracted between the rows and
% both virtual rings are closed by the trace (the boxes above and below).
\newcommand{\TNMPVOverlap}[3]{%
    \begin{tikzpicture}[tn picture]
        \TN@declareWordTensor{ovuL}{(0,0.42)}
        \TN@declareWordTensor{ovuM}{(\TN@chainsep,0.42)}
        \TN@declareWordTensor{ovuR}{(\TN@chainright,0.42)}
        \TN@vconnectports{ovuLE}{ovuMW}
        \TN@vopenport{ovuME}{0.62,0}
        \TN@vopenport{ovuRW}{-0.62,0}
        \TN@vtraceportsabove{ovuLW}{ovuRE}{0.48}{0.32}
        \TN@declareWordTensor{ovdL}{(0,-0.42)}
        \TN@declareWordTensor{ovdM}{(\TN@chainsep,-0.42)}
        \TN@declareWordTensor{ovdR}{(\TN@chainright,-0.42)}
        \TN@vconnectports{ovdLE}{ovdMW}
        \TN@vopenport{ovdME}{0.62,0}
        \TN@vopenport{ovdRW}{-0.62,0}
        \TN@vtraceportsbelow{ovdLW}{ovdRE}{0.48}{0.32}
        \TN@pconnectports{ovuLS}{ovdLN}
        \TN@pconnectports{ovuMS}{ovdMN}
        \TN@pconnectports{ovuRS}{ovdRN}
        % single centred ellipsis in the omitted-site column (no line crosses it)
        \node at (\TN@chainellipsisx,0) {$\cdots$};
        \node[anchor=east] at (-0.85,0.42) {$#1$};
        \node[anchor=east] at (-0.85,-0.42) {$\overline{#2}$};
        \node at (\TN@chainellipsisx,-1.02) {$#3$};
    \end{tikzpicture}%
}

\newcommand{\TNTransferMap}[1]{%
    \begin{tikzpicture}[tn picture]
        \TN@doublelayer{tn}{(0,0.44)}{(0,-0.44)}
        \TN@vopenport{tnUW}{-0.75,0}
        \TN@vopenport{tnDW}{-0.75,0}
        \TN@vopenport{tnUE}{0.75,0}
        \TN@vopenport{tnDE}{0.75,0}
        \node at (-1.18,0) {$X$};
        \node at (1.36,0) {$\E_{#1}(X)$};
        \node at ($(tnUS)!0.5!(tnDN)+(0.24,0)$) {$i$};
    \end{tikzpicture}%
}

\newcommand{\TNMPOCell}[3]{%
    \begin{TNDiagram}[compact]
        \TNMPOSite{tnM}{(0,0)}{#1}
        \TNOpenVirtualWest{tnMWest}
        \TNOpenVirtualEast{tnMEast}
        \TNOpenPhysicalNorth{tnMKet}
        \TNOpenPhysicalSouth{tnMBra}
        \node[above] at (tnMKet) {$#2$};
        \node[below] at (tnMBra) {$#3$};
    \end{TNDiagram}%
}

\newcommand{\TNMPOChain}[6]{%
    \begin{TNDiagram}[normal]
        \TNHorizontalWord{tn}{#2}{#3}{#6}
        \node[above] at (tnRN) {$#4$};
        \node[below] at (tnRS) {$#5$};
        \node[anchor=east] at (tnLW) {$#1$};
    \end{TNDiagram}%
}

% Local purification of an MPO tensor: the ket and bra tensors share the
% ancillary index (straight contraction on the right) and retain their physical
% and virtual legs.
\newcommand{\TNMPOLocalPurification}{%
    \begin{TNEquationRow}[compact]
        \TNTerm{%
            \TNMPOSite{purM}{(0,0)}{M}
            \TNOpenVirtualWest{purMWest}
            \TNOpenVirtualEast{purMEast}
            \TNOpenPhysicalNorth{purMKet}
            \TNOpenPhysicalSouth{purMBra}}
        \TNRelation{=}
        \TNTerm{%
            \TNPurificationSite{pur}{(0,0)}{A}{\overline A}
            \TNOpenVirtualWest{purKetWest}
            \TNOpenVirtualEast{purKetEast}
            \TNOpenVirtualWest{purBraWest}
            \TNOpenVirtualEast{purBraEast}
            \TNOpenPhysicalNorth{purKet}
            \TNOpenPhysicalSouth{purBra}
            \node[right] at (purKetAncilla) {$k$};}
    \end{TNEquationRow}%
}

% The physical renormalization maps of the mixed-state RFP definition.
\newcommand{\TNMPORenormalizationTS}{%
    \begin{TNEquationRow}[compact]
        \ensuremath{M_1(X)=}
        \TNTerm{%
            \TNMPOSite{rtsOne}{(0,0)}{M}
            \TNOpenVirtualWest{rtsOneWest}
            \TNOpenVirtualEast{rtsOneEast}
            \TNOpenPhysicalNorth{rtsOneKet}
            \TNOpenPhysicalSouth{rtsOneBra}}
        \TNRelation{\xrightleftharpoons[\mathcal S]{\mathcal T}}
        \TNTerm{%
            \TNMPOSite{rtsLeft}{(0,0)}{M}
            \TNMPOSite{rtsRight}{(1.25,0)}{M}
            \TNOpenVirtualWest{rtsLeftWest}
            \TNConnectVirtual{rtsLeftEast}{rtsRightWest}
            \TNOpenVirtualEast{rtsRightEast}
            \TNOpenPhysicalNorth{rtsLeftKet}
            \TNOpenPhysicalSouth{rtsLeftBra}
            \TNOpenPhysicalNorth{rtsRightKet}
            \TNOpenPhysicalSouth{rtsRightBra}}
        \ensuremath{=M_2(X)}
    \end{TNEquationRow}%
}

% Contraction of neighbouring right and left sector tensors.
\newcommand{\TNEtaSectorDecomposition}{%
    \begin{TNEquationRow}[compact]
        \ensuremath{\eta_{k,h}=r_k l_h}
    \end{TNEquationRow}%
}

% The two proved two-site operations from Appendix C.2.  The upper row is the
% regrouping followed by T_0's partial trace.  The lower row begins with both
% the outer and neighbouring factors already supplied and draws only the
% inverse regrouping S_2; it deliberately contains no preparation map.
\newcommand{\TNMPDOTwoSiteTraceAndShift}{%
    \begin{tikzpicture}[tn picture, scale=0.72]
        \node[anchor=east] at (-3.72,1.18) {$\mathcal T_0:$};
        \TN@factorwithlegs{tzeroLk}{(-3.15,1.18)}{l_k}
        \TN@factorwithlegs{tzeroRk}{(-2.15,1.18)}{r_k}
        \TN@factorwithlegs{tzeroLh}{(-1.15,1.18)}{l_h}
        \TN@factorwithlegs{tzeroRh}{(-0.15,1.18)}{r_h}
        \TN@vconnect{tzeroLk}{east}{tzeroRk}{west}
        \TN@vconnect{tzeroRk}{east}{tzeroLh}{west}
        \TN@vconnect{tzeroLh}{east}{tzeroRh}{west}
        \TN@grouppath{
            (-2.62,0.68) rectangle (-0.68,1.68);}
        \node at (-1.65,0.44) {$\tr_{R_k\otimes L_h}$};
        \TN@arrow[->, line width=0.55pt]{ (0.52,1.18) -- (1.32,1.18);}
        \TN@factorwithlegs{tzeroOutL}{(1.88,1.18)}{l_k}
        \TN@factorwithlegs{tzeroOutR}{(2.88,1.18)}{r_h}
        \TN@vconnect{tzeroOutL}{east}{tzeroOutR}{west}

        \node[anchor=east] at (-3.72,-1.18) {$\mathcal S_2:$};
        \TN@factorpath{ (-3.42,-1.72) rectangle (-1.38,-0.64);}
        \TN@factorwithlegs{sTwoLk}{(-2.90,-1.18)}{l_k}
        \TN@factorwithlegs{sTwoRh}{(-1.90,-1.18)}{r_h}
        \TN@vconnect{sTwoLk}{east}{sTwoRh}{west}
        \node at (-1.00,-1.18) {$\otimes$};
        \TN@factorpath{ (-0.60,-1.72) rectangle (1.44,-0.64);}
        \TN@factorwithlegs{sTwoRk}{(-0.08,-1.18)}{r_k}
        \TN@factorwithlegs{sTwoLh}{(0.92,-1.18)}{l_h}
        \TN@vconnect{sTwoRk}{east}{sTwoLh}{west}
        \TN@arrow[->, line width=0.55pt]{ (1.72,-1.18) -- (2.52,-1.18);}
        \begin{scope}[xshift=5.98cm]
            \TN@factorwithlegs{sTwoOutLk}{(-3.15,-1.18)}{l_k}
            \TN@factorwithlegs{sTwoOutRk}{(-2.15,-1.18)}{r_k}
            \TN@factorwithlegs{sTwoOutLh}{(-1.15,-1.18)}{l_h}
            \TN@factorwithlegs{sTwoOutRh}{(-0.15,-1.18)}{r_h}
            \TN@vconnect{sTwoOutLk}{east}{sTwoOutRk}{west}
            \TN@vconnect{sTwoOutRk}{east}{sTwoOutLh}{west}
            \TN@vconnect{sTwoOutLh}{east}{sTwoOutRh}{west}
        \end{scope}
        \node at (-1.99,-2.03) {outer factors};
        \node at (0.42,-2.03) {neighbouring factors};
    \end{tikzpicture}%
}

% The three-site analogues.  The dashed middle block is exactly the Omega
% index space traced by S_0.  The lower row starts with that entire block and
% applies only the inverse regrouping T_2, not T_1 or a full recovery map.
\newcommand{\TNMPDOThreeSiteTraceAndShift}{%
    \begin{tikzpicture}[tn picture, scale=0.67]
        \node[anchor=east] at (-4.70,1.35) {$\mathcal S_0:$};
        \TN@factorwithlegs{szeroLk}{(-4.10,1.35)}{l_k}
        \TN@factorwithlegs{szeroRk}{(-3.10,1.35)}{r_k}
        \TN@factorwithlegs{szeroLl}{(-2.10,1.35)}{l_l}
        \TN@factorwithlegs{szeroRl}{(-1.10,1.35)}{r_l}
        \TN@factorwithlegs{szeroLh}{(-0.10,1.35)}{l_h}
        \TN@factorwithlegs{szeroRh}{(0.90,1.35)}{r_h}
        \TN@vconnect{szeroLk}{east}{szeroRk}{west}
        \TN@vconnect{szeroRk}{east}{szeroLl}{west}
        \TN@vconnect{szeroLl}{east}{szeroRl}{west}
        \TN@vconnect{szeroRl}{east}{szeroLh}{west}
        \TN@vconnect{szeroLh}{east}{szeroRh}{west}
        \TN@grouppath{
            (-3.57,0.84) rectangle (0.37,1.86);}
        \node at (-1.60,0.50)
            {$\tr_{\bigoplus_l[(R_k\otimes L_l)\otimes(R_l\otimes L_h)]}$};
        \TN@arrow[->, line width=0.55pt]{ (1.58,1.35) -- (2.38,1.35);}
        \TN@factorwithlegs{szeroOutL}{(2.94,1.35)}{l_k}
        \TN@factorwithlegs{szeroOutR}{(3.94,1.35)}{r_h}
        \TN@vconnect{szeroOutL}{east}{szeroOutR}{west}

        \node[anchor=east] at (-4.70,-1.35) {$\mathcal T_2:$};
        \TN@factorpath{ (-4.37,-1.90) rectangle (-2.33,-0.80);}
        \TN@factorwithlegs{tTwoLk}{(-3.85,-1.35)}{l_k}
        \TN@factorwithlegs{tTwoRh}{(-2.85,-1.35)}{r_h}
        \TN@vconnect{tTwoLk}{east}{tTwoRh}{west}
        \node at (-1.96,-1.35) {$\otimes$};
        \TN@factorpath{ (-1.55,-1.90) rectangle (2.42,-0.80);}
        \TN@factorwithlegs{tTwoRk}{(-1.03,-1.35)}{r_k}
        \TN@factorwithlegs{tTwoLl}{(-0.03,-1.35)}{l_l}
        \TN@factorwithlegs{tTwoRl}{(0.97,-1.35)}{r_l}
        \TN@factorwithlegs{tTwoLh}{(1.97,-1.35)}{l_h}
        \TN@vconnect{tTwoRk}{east}{tTwoLl}{west}
        \TN@vconnect{tTwoLl}{east}{tTwoRl}{west}
        \TN@vconnect{tTwoRl}{east}{tTwoLh}{west}
        \TN@arrow[->, line width=0.55pt]{ (2.68,-1.35) -- (3.48,-1.35);}
        \begin{scope}[xshift=7.90cm]
            \TN@factorwithlegs{tTwoOutLk}{(-4.10,-1.35)}{l_k}
            \TN@factorwithlegs{tTwoOutRk}{(-3.10,-1.35)}{r_k}
            \TN@factorwithlegs{tTwoOutLl}{(-2.10,-1.35)}{l_l}
            \TN@factorwithlegs{tTwoOutRl}{(-1.10,-1.35)}{r_l}
            \TN@factorwithlegs{tTwoOutLh}{(-0.10,-1.35)}{l_h}
            \TN@factorwithlegs{tTwoOutRh}{(0.90,-1.35)}{r_h}
            \TN@vconnect{tTwoOutLk}{east}{tTwoOutRk}{west}
            \TN@vconnect{tTwoOutRk}{east}{tTwoOutLl}{west}
            \TN@vconnect{tTwoOutLl}{east}{tTwoOutRl}{west}
            \TN@vconnect{tTwoOutRl}{east}{tTwoOutLh}{west}
            \TN@vconnect{tTwoOutLh}{east}{tTwoOutRh}{west}
        \end{scope}
        \node at (-3.35,-2.24) {outer factors};
        \node at (0.44,-2.24) {$\Omega_{k,h}=\bigoplus_l(\cdots)$};
    \end{tikzpicture}%
}

% The three stages of the physical-sector refinement channel in Appendix C.2.
% The layout follows MPDO_TMatrix.png: a two-site sector row, the boundary row
% left by T_0, and the three-site factor order obtained after T_1 and T_2.
\newcommand{\TNMPDORefinementConstruction}{%
    \begin{tikzpicture}[tn picture, scale=0.85]
        \node[anchor=east] at (-5.35,2.45) {$\mathfrak R_2:$};
        \node[anchor=east] at (-4.30,2.45) {$\displaystyle\bigoplus_{k,h}$};
        \TN@factorwithlegs{trefTopLk}{(-3.72,2.45)}{l_k}
        \TN@factorwithlegs{trefTopRk}{(-2.72,2.45)}{r_k}
        \TN@factorwithlegs{trefTopLh}{(-1.72,2.45)}{l_h}
        \TN@factorwithlegs{trefTopRh}{(-0.72,2.45)}{r_h}
        \TN@vconnect{trefTopLk}{east}{trefTopRk}{west}
        \TN@vconnect{trefTopRk}{east}{trefTopLh}{west}
        \TN@vconnect{trefTopLh}{east}{trefTopRh}{west}
        \TN@grouppath{
            (-3.25,1.94) rectangle (-1.19,2.96);}
        \node at (-2.22,1.68) {$\tr_{R_k\otimes L_h}$};

        \TN@arrow[->, line width=0.55pt]{ (2.15,2.18) -- (2.15,1.22)
            node[midway, right] {$\mathcal T_0$};}
        \node[anchor=east] at (-2.08,0.86)
            {$\displaystyle\bigoplus_{k,h}a_kb_h$};
        \TN@factorwithlegs{trefMidLk}{(-1.50,0.86)}{l_k}
        \TN@factorwithlegs{trefMidRh}{(-0.50,0.86)}{r_h}
        \TN@vconnect{trefMidLk}{east}{trefMidRh}{west}

        \TN@arrow[->, line width=0.55pt]{ (2.15,0.38) -- (2.15,-0.58)
            node[midway, right] {$\mathcal T_2\mathcal T_1$};}
        \node[anchor=east] at (-5.35,-0.96) {$\mathfrak R_3:$};
        \node[anchor=east] at (-4.30,-0.96)
            {$\displaystyle\bigoplus_{k,h,l}$};
        \TN@factorwithlegs{trefBotLk}{(-3.72,-0.96)}{l_k}
        \TN@factorwithlegs{trefBotRk}{(-2.72,-0.96)}{r_k}
        \TN@factorwithlegs{trefBotLl}{(-1.72,-0.96)}{l_l}
        \TN@factorwithlegs{trefBotRl}{(-0.72,-0.96)}{r_l}
        \TN@factorwithlegs{trefBotLh}{(0.28,-0.96)}{l_h}
        \TN@factorwithlegs{trefBotRh}{(1.28,-0.96)}{r_h}
        \TN@vconnect{trefBotLk}{east}{trefBotRk}{west}
        \TN@vconnect{trefBotRk}{east}{trefBotLl}{west}
        \TN@vconnect{trefBotLl}{east}{trefBotRl}{west}
        \TN@vconnect{trefBotRl}{east}{trefBotLh}{west}
        \TN@vconnect{trefBotLh}{east}{trefBotRh}{west}
    \end{tikzpicture}%
}

% The refinement half of MPDO_TSKKK.png.  This records only the direction and
% site counts of the channel; the closure identity belongs to its theorem.
\newcommand{\TNMPDORefinementDirection}{%
    \begin{tikzpicture}[tn picture, scale=0.85]
        \TN@factorwithlegs{trefDirTwoL}{(-2.85,0)}{\mathcal K}
        \TN@factorwithlegs{trefDirTwoR}{(-1.75,0)}{\mathcal K}
        \TN@vconnect{trefDirTwoL}{east}{trefDirTwoR}{west}
        \TN@vpath{ (trefDirTwoL.west) -- ++(-0.48,0);}
        \TN@vpath{ (trefDirTwoR.east) -- ++(0.48,0);}
        \node at (-2.30,-0.92) {$\mathfrak R_2({\cal K}_2(X))$};

        \TN@arrow[->, line width=0.65pt, bend left=18]{
            (-0.94,0.28) to node[above] {$\mathcal T$} (0.04,0.28);}

        \TN@factorwithlegs{trefDirThreeL}{(0.85,0)}{\mathcal K}
        \TN@factorwithlegs{trefDirThreeM}{(1.95,0)}{\mathcal K}
        \TN@factorwithlegs{trefDirThreeR}{(3.05,0)}{\mathcal K}
        \TN@vconnect{trefDirThreeL}{east}{trefDirThreeM}{west}
        \TN@vconnect{trefDirThreeM}{east}{trefDirThreeR}{west}
        \TN@vpath{ (trefDirThreeL.west) -- ++(-0.48,0);}
        \TN@vpath{ (trefDirThreeR.east) -- ++(0.48,0);}
        \node at (1.95,-0.92) {$\mathfrak R_3({\cal K}_3(X))$};
    \end{tikzpicture}%
}

% The overlapping applications of T and S used after two-site blocking in
% Theorem 4.9.  The endpoint frames identify one and two blocked sites; the
% intermediate rows retain the original physical-site decomposition.
\newcommand{\TNMPDOBlockedRFPChannels}{%
    \begin{tikzpicture}[tn picture, scale=0.72]
        \foreach \x/\name in {-5.0/btTwoA,-4.0/btTwoB} {
            \TN@factorwithlegs{\name}{(\x,1.35)}{\widehat{\cal K}}
        }
        \TN@vconnect{btTwoA}{east}{btTwoB}{west}
        \TN@grouppath{
            (-5.48,0.87) rectangle (-3.52,1.83);}
        \node at (-4.50,0.54) {one blocked site};

        \foreach \x/\name in {-1.4/btThreeA,-0.4/btThreeB,0.6/btThreeC} {
            \TN@factorwithlegs{\name}{(\x,1.35)}{\widehat{\cal K}}
        }
        \TN@vconnect{btThreeA}{east}{btThreeB}{west}
        \TN@vconnect{btThreeB}{east}{btThreeC}{west}

        \foreach \x/\name in {3.0/btFourA,4.0/btFourB,5.0/btFourC,6.0/btFourD} {
            \TN@factorwithlegs{\name}{(\x,1.35)}{\widehat{\cal K}}
        }
        \TN@vconnect{btFourA}{east}{btFourB}{west}
        \TN@vconnect{btFourB}{east}{btFourC}{west}
        \TN@vconnect{btFourC}{east}{btFourD}{west}
        \TN@grouppath{
            (2.52,0.87) rectangle (4.48,1.83);}
        \TN@grouppath{
            (4.52,0.87) rectangle (6.48,1.83);}
        \node at (4.50,0.54) {two blocked sites};

        \TN@arrow[->, line width=0.60pt]{ (-3.22,1.35) -- (-1.88,1.35)
            node[midway, above] {$\widehat{\mathcal T}$};}
        \TN@arrow[->, line width=0.60pt]{ (1.08,1.35) -- (2.52,1.35)
            node[midway, above] {$\widehat{\mathcal T}_{(12)3}$};}
        \node at (0.50,2.25)
            {$\widehat{\mathcal T}^2:\widehat{\cal K}_2(X)
              \longmapsto\widehat{\cal K}_4(X)$};

        \foreach \x/\name in {3.0/bsFourA,4.0/bsFourB,5.0/bsFourC,6.0/bsFourD} {
            \TN@factorwithlegs{\name}{(\x,-1.30)}{\widehat{\cal K}}
        }
        \TN@vconnect{bsFourA}{east}{bsFourB}{west}
        \TN@vconnect{bsFourB}{east}{bsFourC}{west}
        \TN@vconnect{bsFourC}{east}{bsFourD}{west}
        \TN@grouppath{
            (2.52,-1.78) rectangle (4.48,-0.82);}
        \TN@grouppath{
            (4.52,-1.78) rectangle (6.48,-0.82);}
        \foreach \x/\name in {-1.4/bsThreeA,-0.4/bsThreeB,0.6/bsThreeC} {
            \TN@factorwithlegs{\name}{(\x,-1.30)}{\widehat{\cal K}}
        }
        \TN@vconnect{bsThreeA}{east}{bsThreeB}{west}
        \TN@vconnect{bsThreeB}{east}{bsThreeC}{west}
        \foreach \x/\name in {-5.0/bsTwoA,-4.0/bsTwoB} {
            \TN@factorwithlegs{\name}{(\x,-1.30)}{\widehat{\cal K}}
        }
        \TN@vconnect{bsTwoA}{east}{bsTwoB}{west}
        \TN@grouppath{
            (-5.48,-1.78) rectangle (-3.52,-0.82);}
        \TN@arrow[<-, line width=0.60pt]{ (1.08,-1.30) -- (2.52,-1.30)
            node[midway, above] {$\widehat{\mathcal S}_{(123)4}$};}
        \TN@arrow[<-, line width=0.60pt]{ (-3.22,-1.30) -- (-1.88,-1.30)
            node[midway, above] {$\widehat{\mathcal S}$};}
        \node at (0.50,-2.20)
            {$\widehat{\mathcal S}^2:\widehat{\cal K}_4(X)
              \longmapsto\widehat{\cal K}_2(X)$};
    \end{tikzpicture}%
}

% Transport of the two-step channels from the sector coordinates selected by
% the physical isometry back to the original physical coordinates. The
% horizontal arrows follow MPDO_TSKKK.png; the vertical conjugations follow
% the U and U^dagger convention of MPDO_K=sum_k_LkRk.png.
\newcommand{\TNMPDOPhysicalIsometryTransport}{%
    \begin{tikzpicture}[tn picture, scale=0.64]
        \foreach \x/\name in {-5.0/ptRawTwoA,-4.0/ptRawTwoB} {
            \TN@factorwithlegs{\name}{(\x,1.45)}{\cal K}
        }
        \TN@vconnect{ptRawTwoA}{east}{ptRawTwoB}{west}
        \TN@grouppath{
            (-5.48,0.97) rectangle (-3.52,1.93);}
        \node at (-4.50,2.27) {${\cal K}_2(X)$};

        \foreach \x/\name in {3.0/ptRawFourA,4.0/ptRawFourB,
                5.0/ptRawFourC,6.0/ptRawFourD} {
            \TN@factorwithlegs{\name}{(\x,1.45)}{\cal K}
        }
        \TN@vconnect{ptRawFourA}{east}{ptRawFourB}{west}
        \TN@vconnect{ptRawFourB}{east}{ptRawFourC}{west}
        \TN@vconnect{ptRawFourC}{east}{ptRawFourD}{west}
        \TN@grouppath{
            (2.52,0.97) rectangle (4.48,1.93);}
        \TN@grouppath{
            (4.52,0.97) rectangle (6.48,1.93);}
        \node at (4.50,2.27) {${\cal K}_4(X)$};

        \foreach \x/\name in {-5.0/ptHatTwoA,-4.0/ptHatTwoB} {
            \TN@factorwithlegs{\name}{(\x,-1.45)}{\widehat{\cal K}}
        }
        \TN@vconnect{ptHatTwoA}{east}{ptHatTwoB}{west}
        \TN@grouppath{
            (-5.48,-1.93) rectangle (-3.52,-0.97);}
        \node at (-4.50,-2.27) {$\widehat{\cal K}_2(X)$};

        \foreach \x/\name in {3.0/ptHatFourA,4.0/ptHatFourB,
                5.0/ptHatFourC,6.0/ptHatFourD} {
            \TN@factorwithlegs{\name}{(\x,-1.45)}{\widehat{\cal K}}
        }
        \TN@vconnect{ptHatFourA}{east}{ptHatFourB}{west}
        \TN@vconnect{ptHatFourB}{east}{ptHatFourC}{west}
        \TN@vconnect{ptHatFourC}{east}{ptHatFourD}{west}
        \TN@grouppath{
            (2.52,-1.93) rectangle (4.48,-0.97);}
        \TN@grouppath{
            (4.52,-1.93) rectangle (6.48,-0.97);}
        \node at (4.50,-2.27) {$\widehat{\cal K}_4(X)$};

        \TN@arrow[->, line width=0.60pt, bend left=13]{
            (-3.22,1.62) to node[above] {$\mathcal T_{\cal K}^2$} (2.22,1.62);}
        \TN@arrow[->, line width=0.60pt, bend left=13]{
            (2.22,1.28) to node[below] {$\mathcal S_{\cal K}^2$} (-3.22,1.28);}
        \TN@arrow[->, line width=0.60pt, bend left=13]{
            (-3.22,-1.28) to node[above] {$\widehat{\mathcal T}^2$} (2.22,-1.28);}
        \TN@arrow[->, line width=0.60pt, bend left=13]{
            (2.22,-1.62) to node[below] {$\widehat{\mathcal S}^2$} (-3.22,-1.62);}

        \TN@arrow[->, line width=0.55pt]{ (-4.64,0.68) -- (-4.64,-0.68)
            node[midway, left] {$U^{\otimes 2}$};}
        \TN@arrow[->, line width=0.55pt]{ (-4.36,-0.68) -- (-4.36,0.68)
            node[midway, right] {$(U^\dagger)^{\otimes 2}$};}
        \TN@arrow[->, line width=0.55pt]{ (4.36,0.68) -- (4.36,-0.68)
            node[midway, left] {$U^{\otimes 4}$};}
        \TN@arrow[->, line width=0.55pt]{ (4.64,-0.68) -- (4.64,0.68)
            node[midway, right] {$(U^\dagger)^{\otimes 4}$};}
    \end{tikzpicture}%
}

% Fixed-sector factorization of the arbitrary-X two-site physical closure.
% The left-hand row shows the sector subspins and the outer virtual closure;
% the right-hand side shows only the proved boundary/neighbor tensor product.
\newcommand{\TNMPDOTwoSiteClosureFactorization}{%
    \begin{tikzpicture}[tn picture, scale=0.50]
        \TN@factorwithlegs{ctwoLk}{(-4.40,0.55)}{l_k}
        \TN@factorwithlegs{ctwoRk}{(-3.40,0.55)}{r_k}
        \TN@factorwithlegs{ctwoLh}{(-2.40,0.55)}{l_h}
        \TN@factorwithlegs{ctwoRh}{(-1.40,0.55)}{r_h}
        \TN@vconnect{ctwoLk}{east}{ctwoRk}{west}
        \TN@vconnect{ctwoRk}{east}{ctwoLh}{west}
        \TN@vconnect{ctwoLh}{east}{ctwoRh}{west}
        \TN@insertion[label=below:$X$]{ctwoX}{(-2.90,-0.95)}
        \TN@vtracepath{ (ctwoLk.west)
          .. controls (-5.05,0.55) and (-5.05,-0.95) .. (ctwoX.west);}
        \TN@vtracepath{ (ctwoX.east)
          .. controls (-0.75,-0.95) and (-0.75,0.55) .. (ctwoRh.east);}
        \node at (0.00,0.05) {$=$};
        \TN@component[minimum width=2.25cm, minimum height=0.72cm]
            {ctwoB}{(1.70,0.05)}{B_{k,h}(X)}
        \node at (3.20,0.05) {$\otimes$};
        \TN@factor[minimum width=1.75cm, minimum height=0.72cm]
            {ctwoEta}{(4.45,0.05)}{\eta_{k,h}}
        \node at (1.70,-0.62) {$L_k\otimes R_h$};
        \node at (4.45,-0.62) {$R_k\otimes L_h$};
    \end{tikzpicture}%
}

% Fixed-sector factorization of the arbitrary-X three-site physical closure.
% No preparation, positivity, or channel identity is represented.
\newcommand{\TNMPDOThreeSiteClosureFactorization}{%
    \begin{tikzpicture}[tn picture, scale=0.43]
        \TN@factorwithlegs{cthreeLk}{(-5.70,0.62)}{l_k}
        \TN@factorwithlegs{cthreeRk}{(-4.70,0.62)}{r_k}
        \TN@factorwithlegs{cthreeLl}{(-3.70,0.62)}{l_l}
        \TN@factorwithlegs{cthreeRl}{(-2.70,0.62)}{r_l}
        \TN@factorwithlegs{cthreeLh}{(-1.70,0.62)}{l_h}
        \TN@factorwithlegs{cthreeRh}{(-0.70,0.62)}{r_h}
        \TN@vconnect{cthreeLk}{east}{cthreeRk}{west}
        \TN@vconnect{cthreeRk}{east}{cthreeLl}{west}
        \TN@vconnect{cthreeLl}{east}{cthreeRl}{west}
        \TN@vconnect{cthreeRl}{east}{cthreeLh}{west}
        \TN@vconnect{cthreeLh}{east}{cthreeRh}{west}
        \TN@insertion[label=below:$X$]{cthreeX}{(-3.20,-1.00)}
        \TN@vtracepath{ (cthreeLk.west)
          .. controls (-6.35,0.62) and (-6.35,-1.00) .. (cthreeX.west);}
        \TN@vtracepath{ (cthreeX.east)
          .. controls (-0.05,-1.00) and (-0.05,0.62) .. (cthreeRh.east);}
        \node at (0.45,0.05) {$=$};
        \TN@component[minimum width=2.05cm, minimum height=0.72cm]
            {cthreeB}{(1.85,0.05)}{B_{k,h}(X)}
        \node at (3.15,0.05) {$\otimes$};
        \TN@factor[minimum width=1.65cm, minimum height=0.72cm]
            {cthreeEtaKL}{(4.20,0.05)}{\eta_{k,l}}
        \node at (5.25,0.05) {$\otimes$};
        \TN@factor[minimum width=1.65cm, minimum height=0.72cm]
            {cthreeEtaLH}{(6.30,0.05)}{\eta_{l,h}}
        \node at (1.85,-0.65) {$L_k\otimes R_h$};
        \node at (5.25,-0.65)
          {$(R_k\otimes L_l)\otimes(R_l\otimes L_h)$};
    \end{tikzpicture}%
}

% Adjacent-bond commutativity on one fixed sector triple.  The six vertical
% lines are the factors L_k, R_k, L_l, R_l, L_h, R_h.  Reading from bottom to
% top, the two panels differ only by the order of the operators on the disjoint
% pairs R_k \otimes L_l and R_l \otimes L_h.
\newcommand{\TNMPDOFixedSectorAdjacentCommutativity}{%
    \begin{tikzpicture}[tn picture, scale=0.78]
        % The four outer factors do not meet an operator in this comparison.
        \foreach \x in {-5.25,-1.25,1.25,5.25} {
            \TN@vpath{ (\x,-1.12) -- (\x,1.12);}
        }

        \TN@factor[minimum width=1.38cm, minimum height=0.48cm,
            inner sep=1pt]{commEtaKLLeft}{(-4.05,0.38)}{\eta_{k,l}}
        \TN@factor[minimum width=1.38cm, minimum height=0.48cm,
            inner sep=1pt]{commEtaLHLeft}{(-2.45,-0.38)}{\eta_{l,h}}

        \TN@factor[minimum width=1.38cm, minimum height=0.48cm,
            inner sep=1pt]{commEtaKLRight}{(2.45,-0.38)}{\eta_{k,l}}
        \TN@factor[minimum width=1.38cm, minimum height=0.48cm,
            inner sep=1pt]{commEtaLHRight}{(4.05,0.38)}{\eta_{l,h}}

        % Each operator carries two factors.  The eight affected index lines
        % terminate on named boundary ports rather than passing underneath the
        % filled factor plates.
        \foreach \factor/\low/\high in {
            commEtaKLLeft/-1.26/0.50,
            commEtaLHLeft/-0.50/1.26,
            commEtaKLRight/-0.50/1.26,
            commEtaLHRight/-1.26/0.50} {
            \TN@vsouthport{\factor SouthLeft}{\factor}{0.21}
            \TN@vsouthport{\factor SouthRight}{\factor}{0.79}
            \TN@vnorthport{\factor NorthLeft}{\factor}{0.21}
            \TN@vnorthport{\factor NorthRight}{\factor}{0.79}
            \TN@vopenport{\factor SouthLeft}{0,\low}
            \TN@vopenport{\factor SouthRight}{0,\low}
            \TN@vopenport{\factor NorthLeft}{0,\high}
            \TN@vopenport{\factor NorthRight}{0,\high}
        }

        \node at (-3.25,1.57) {$C_{01}^{k,l,h}C_{12}^{k,l,h}$};
        \node at (3.25,1.57) {$C_{12}^{k,l,h}C_{01}^{k,l,h}$};
        \node at (0,0) {$=$};

        \foreach \x/\lab in {
            -5.25/{B_k^L},-4.45/{B_k^R},-3.65/{B_l^L},-2.85/{B_l^R},
            -2.05/{B_h^L},-1.25/{B_h^R},
             1.25/{B_k^L}, 2.05/{B_k^R}, 2.85/{B_l^L}, 3.65/{B_l^R},
             4.45/{B_h^L}, 5.25/{B_h^R}} {
            \node at (\x,-1.42) {$\lab$};
        }
        \foreach \a/\b/\lab in {
            -5.58/-4.12/0,-3.98/-2.52/1,-2.38/-0.92/2,
             0.92/2.38/0, 2.52/3.98/1, 4.12/5.58/2} {
            \TN@vpath{ (\a,-1.72) -- (\a,-1.84) -- (\b,-1.84) -- (\b,-1.72);}
            \node at ({(\a+\b)/2},-2.05) {site $\lab$};
        }
    \end{tikzpicture}%
}

% A closed chain of factorized sector tensors regrouped into neighbouring
% eta operators.  This mirrors the cyclic contractions in Appendix C.2.
\newcommand{\TNMPDOCyclicEtaContraction}{%
    \begin{tikzpicture}[tn picture]
        \foreach \name/\x/\lab in {Zero/-2.70/0,One/-0.90/1,Last/1.30/{N-1}} {
            \TN@factor[minimum width=0.58cm, minimum height=0.38cm,
                inner sep=1pt]{tnL\name}{(\x,0.34)}{l_{k_\lab}}
            \TN@factor[minimum width=0.58cm, minimum height=0.38cm,
                inner sep=1pt]{tnR\name}{(\x+1.14,0.34)}{r_{k_\lab}}
            \node at (\x+0.57,0.34) {$\otimes$};
            \TN@pleg{tnL\name}{north}{0,0.42}
            \TN@pleg{tnL\name}{south}{0,-0.42}
            \TN@pleg{tnR\name}{north}{0,0.42}
            \TN@pleg{tnR\name}{south}{0,-0.42}
        }
        \TN@vpath{ (tnRZero.east) -- (tnLOne.west);}
        \TN@vpath{ (tnROne.east) -- (0.62,0.34);}
        \TN@vpath{ (0.94,0.34) -- (tnLLast.west);}
        \node at (0.78,0.34) {$\cdots$};
        \TN@vtracepath{ (tnRLast.east)
            .. controls (3.04,0.34) and (3.04,-0.42) .. (2.60,-0.42)
            -- (-3.04,-0.42)
            .. controls (-3.46,-0.42) and (-3.46,0.34) .. (tnLZero.west);}
        \node at (0,-0.90)
            {$\eta_{k_0,k_1}\otimes\eta_{k_1,k_2}\otimes\cdots
              \otimes\eta_{k_{N-1},k_0}$};
    \end{tikzpicture}%
}

% Applying the inverse tensor to the first and third sites of a three-site
% MPO word leaves one matrix entry of the middle tensor and one entry of the
% complementary virtual tail.
\newcommand{\TNMPDOInverseMapThreeSiteContraction}{%
    \begin{tikzpicture}[tn picture]
        % Three-site MPO word.  At the first and third sites, both physical
        % legs are contracted against the corresponding inverse tensor.
        \TNTensor{tnKone}{(-3.00,0)}
        \TNTensor{tnKtwo}{(-1.60,0)}
        \TNTensor{tnKthree}{(-0.20,0)}
        \TN@vpath{ (tnKone.east) -- (tnKtwo.west);}
        \TN@vpath{ (tnKtwo.east) -- (tnKthree.west);}
        \node[anchor=east] at (-3.14,-0.28) {${\cal K}$};
        \node[anchor=east] at (-1.74,-0.28) {${\cal K}$};
        \node[anchor=east] at (-0.34,-0.28) {${\cal K}$};

        \TNTensor{tnInvOne}{(-3.00,1.18)}
        \TNTensor{tnInvThree}{(-0.20,1.18)}
        \TN@ppath{ (tnKone.north) -- (tnInvOne.south);}
        \TN@ppath{ (tnKthree.north) -- (tnInvThree.south);}
        \TN@ppath{
            (tnKone.south) .. controls (-2.18,-0.38) and (-2.18,1.62) ..
            (tnInvOne.north);}
        \TN@ppath{
            (tnKthree.south) .. controls (-0.98,-0.38) and (-0.98,1.62) ..
            (tnInvThree.north);}
        \TN@vleg{tnInvOne}{west}{-0.52,0}
        \TN@vleg{tnInvOne}{east}{0.52,0}
        \TN@vleg{tnInvThree}{west}{-0.52,0}
        \TN@vleg{tnInvThree}{east}{0.52,0}
        \node[anchor=east] at (-3.12,1.55) {${\cal K}^{-1}$};
        \node[anchor=west] at (-0.08,1.55) {${\cal K}^{-1}$};
        \node[anchor=east] at (-3.58,1.18) {$\alpha_1$};
        \node[anchor=west] at (-2.42,1.18) {$\beta_1$};
        \node[anchor=east] at (-0.78,1.18) {$\alpha_3$};
        \node[anchor=west] at (0.38,1.18) {$\beta_3$};
        \TN@pleg{tnKtwo}{north}{0,0.52}
        \TN@pleg{tnKtwo}{south}{0,-0.52}
        \node at (-1.60,0.82) {$i_2$};
        \node at (-1.60,-0.82) {$j_2$};
        \node at (-2.88,0.58) {$i_1$};
        \node at (-2.08,0.48) {$j_1$};
        \node at (-0.08,0.58) {$i_3$};
        \node at (-1.08,0.48) {$j_3$};

        % The trace closes the three-site word through the remaining sites R.
        \TNComponent{tnTail}{(-1.60,-1.42)}{R}
        \TN@vtracepath{ (tnKone.west) -- ++(-0.42,0)
            |- ($(tnTail.east)+(0,-0.40)$) -- (tnTail.east);}
        \TN@vtracepath{ (tnKthree.east) -- ++(0.42,0)
            |- ($(tnTail.west)+(0,0.38)$) -- (tnTail.west);}

        \node at (0.88,0.20) {$=$};

        % The two surviving factors.
        \TNTensor{tnMiddle}{(2.25,0.20)}
        \TN@vleg{tnMiddle}{west}{-0.58,0}
        \TN@vleg{tnMiddle}{east}{0.58,0}
        \TN@pleg{tnMiddle}{north}{0,0.52}
        \TN@pleg{tnMiddle}{south}{0,-0.52}
        \node[anchor=east] at (2.11,-0.08) {${\cal K}$};
        \node at (2.25,1.02) {$i_2$};
        \node at (2.25,-0.62) {$j_2$};
        \node[anchor=east] at (1.61,0.20) {$\beta_1$};
        \node[anchor=west] at (2.89,0.20) {$\alpha_3$};
        \node at (3.50,0.20) {$\times$};
        \TNComponent{tnTailEntry}{(4.65,0.20)}{R}
        \TN@vpath{ (tnTailEntry.west) -- ++(-0.58,0);}
        \TN@vpath{ (tnTailEntry.east) -- ++(0.58,0);}
        \node[anchor=east] at (4.01,0.20) {$\beta_3$};
        \node[anchor=west] at (5.29,0.20) {$\alpha_1$};
    \end{tikzpicture}%
}

% Horizontal MPDO operator: one row of MPO sites with the virtual bond
% contracted along the row and closed by the trace; the ket and bra legs
% stay external.
\newcommand{\TNMPDOHorizontalOperator}{%
    \begin{tikzpicture}[tn picture]
        \node[anchor=east] at (-0.72,0) {$H^{(N)}(\tilde M)\;=$};
        \TN@closedMPOword{hop}
    \end{tikzpicture}%
}

% The three first-site contractions of the horizontal MPDO operator: the
% matrix P inserted on the ket leg, on the bra leg, and on both legs of the
% first site.
\newcommand{\TNMPDOFirstSiteContractions}{%
    \begin{tikzpicture}[tn picture]
        \begin{scope}
            \TN@closedMPOword{fsl}
            \TN@insertion[label=left:$P$]{fslP}{(0,0.26)}
            \node[anchor=west] at (3.80,0) {$PH^{(N+1)}$};
        \end{scope}
        \begin{scope}[yshift=-2.35cm]
            \TN@closedMPOword{fsr}
            \TN@insertion[label=left:$P$]{fsrP}{(0,-0.26)}
            \node[anchor=west] at (3.80,0) {$H^{(N+1)}P$};
        \end{scope}
        \begin{scope}[yshift=-4.70cm]
            \TN@closedMPOword{fsc}
            \TN@insertion[label=left:$P$]{fscPu}{(0,0.26)}
            \TN@insertion[label=left:$P$]{fscPd}{(0,-0.26)}
            \node[anchor=west] at (3.80,0) {$PH^{(N+1)}P$};
        \end{scope}
    \end{tikzpicture}%
}

% Vertical direct-sum form of the doubled tensor: after the physical isometry
% each block factorizes into a diagonal weight on the multiplicity factor and
% a basis tensor carrying the virtual legs.
\newcommand{\TNMPDOVerticalDirectSum}{%
    \begin{tikzpicture}[tn picture]
        \TNMap{vdUu}{(-2.55,0.76)}{U^\dagger}
        \TNTensor{vdM}{(-2.55,0)}
        \TNMap{vdUd}{(-2.55,-0.76)}{U}
        \TN@pnorthport{vdUuNLeft}{vdUu}{0.30}
        \TN@pnorthport{vdUuNRight}{vdUu}{0.70}
        \TN@psouthport{vdUdSLeft}{vdUd}{0.30}
        \TN@psouthport{vdUdSRight}{vdUd}{0.70}
        \TN@popenport{vdUuNLeft}{0,0.36}
        \TN@popenport{vdUuNRight}{0,0.36}
        \TN@ppath{ (vdUu.south) -- (vdM.north);}
        \TN@ppath{ (vdM.south) -- (vdUd.north);}
        \TN@popenport{vdUdSLeft}{0,-0.36}
        \TN@popenport{vdUdSRight}{0,-0.36}
        \TN@vleg{vdM}{west}{-0.58,0}
        \TN@vleg{vdM}{east}{0.58,0}
        \node at (-2.92,-0.26) {$\tilde M$};
        \node at (-1.43,0) {$=$};
        \node at (-0.62,0) {$\displaystyle\bigoplus_\alpha$};
        \TN@insertion{vdMu}{(0.45,0)}
        \TN@ppath{ (vdMu.north) -- ++(0,0.40);}
        \TN@ppath{ (vdMu.south) -- ++(0,-0.40);}
        \node at (0.45,0.74) {$\mu_\alpha$};
        \TNTensor{vdMa}{(1.55,0)}
        \TN@pleg{vdMa}{north}{0,0.44}
        \TN@pleg{vdMa}{south}{0,-0.44}
        \TN@vleg{vdMa}{west}{-0.55,0}
        \TN@vleg{vdMa}{east}{0.55,0}
        \node at (1.55,0.74) {$M_\alpha$};
    \end{tikzpicture}%
}

% Complete orthogonal reducing-sector decomposition of the vertical tensor.
% The physical isometries V_k surround the irreducible corner B_k; their
% range projections resolve the physical identity and reduce every letter.
\newcommand{\TNMPDOVerticalReducingSectors}{%
    \begin{tikzpicture}[tn picture, scale=0.94]
        \path[draw=black, opacity=0, use as bounding box]
            (-4.35,-1.42) rectangle (5.75,1.42);
        \TNTensor{vrsM}{(-3.50,0)}
        \TN@vleg{vrsM}{west}{-0.58,0}
        \TN@vleg{vrsM}{east}{0.58,0}
        \TN@pleg{vrsM}{north}{0,0.48}
        \TN@pleg{vrsM}{south}{0,-0.48}
        \node[anchor=west] at (-3.33,0.30) {$\widetilde M$};
        \node at (-2.45,0) {$=$};
        \node at (-1.55,0) {$\displaystyle\sum_k$};
        \TNMap{vrsV}{(0,0.82)}{V_k}
        \TNTensor{vrsB}{(0,0)}
        \TNMap{vrsVd}{(0,-0.82)}{V_k^\dagger}
        \TN@ppath{ (vrsV.north) -- ++(0,0.34);}
        \TN@ppath{ (vrsV.south) -- (vrsB.north);}
        \TN@ppath{ (vrsB.south) -- (vrsVd.north);}
        \TN@ppath{ (vrsVd.south) -- ++(0,-0.34);}
        \TN@vleg{vrsB}{west}{-0.60,0}
        \TN@vleg{vrsB}{east}{0.60,0}
        \node[anchor=west] at (0.16,0.27) {$B_k$};
        \node[anchor=west, align=left] at (1.62,0.48)
            {$P_k=V_kV_k^\dagger,\qquad \sum_kP_k=I,$};
        \node[anchor=west, align=left] at (1.62,-0.18)
            {$P_kP_l=0\ (k\ne l),$};
        \node[anchor=west, align=left] at (1.62,-0.78)
            {$P_k\widetilde M=\widetilde MP_k.$};
    \end{tikzpicture}%
}

% Equality of the Gram conjugations for two gauge-equivalent copies of one
% vertical normal tensor. The ordering of the four gauge factors follows
% DAVID_Fig8.png in the proof of Proposition 4.13.
\newcommand{\TNMPDOVerticalGaugeGramComparison}{%
    \begin{tikzpicture}[tn picture]
        \path[draw=black, opacity=0, use as bounding box]
            (-3.45,-1.62) rectangle (3.45,1.62);
        \begin{scope}[xshift=-1.80cm]
            \TNTensor{vggL}{(0,0)}
            \TN@ppath{ (vggL.west) -- ++(-0.72,0);}
            \TN@ppath{ (vggL.east) -- ++(0.72,0);}
            \TN@insertion[label=right:$X_{\alpha,k}$]{vggLX}{(0,0.58)}
            \TN@insertion[label=right:$X_{\alpha,k}^{\dagger}$]{vggLXd}{(0,1.18)}
            \TN@insertion[label=right:$X_{\alpha,k}^{-1}$]{vggLXi}{(0,-0.58)}
            \TN@insertion[label=right:$(X_{\alpha,k}^{-1})^{\dagger}$]{vggLXid}{(0,-1.18)}
            \TN@vpath{ (vggL.north) -- (vggLX.south);}
            \TN@vpath{ (vggLX.north) -- (vggLXd.south);}
            \TN@vpath{ (vggLXd.north) -- ++(0,0.30);}
            \TN@vpath{ (vggL.south) -- (vggLXi.north);}
            \TN@vpath{ (vggLXi.south) -- (vggLXid.north);}
            \TN@vpath{ (vggLXid.south) -- ++(0,-0.30);}
            \node[anchor=east] at (-0.16,0.18) {$M_\alpha$};
        \end{scope}
        \node at (0,0) {$=$};
        \begin{scope}[xshift=1.80cm]
            \TNTensor{vggR}{(0,0)}
            \TN@ppath{ (vggR.west) -- ++(-0.72,0);}
            \TN@ppath{ (vggR.east) -- ++(0.72,0);}
            \TN@insertion[label=right:$X_{\alpha,1}$]{vggRX}{(0,0.58)}
            \TN@insertion[label=right:$X_{\alpha,1}^{\dagger}$]{vggRXd}{(0,1.18)}
            \TN@insertion[label=right:$X_{\alpha,1}^{-1}$]{vggRXi}{(0,-0.58)}
            \TN@insertion[label=right:$(X_{\alpha,1}^{-1})^{\dagger}$]{vggRXid}{(0,-1.18)}
            \TN@vpath{ (vggR.north) -- (vggRX.south);}
            \TN@vpath{ (vggRX.north) -- (vggRXd.south);}
            \TN@vpath{ (vggRXd.north) -- ++(0,0.30);}
            \TN@vpath{ (vggR.south) -- (vggRXi.north);}
            \TN@vpath{ (vggRXi.south) -- (vggRXid.north);}
            \TN@vpath{ (vggRXid.south) -- ++(0,-0.30);}
            \node[anchor=east] at (-0.16,0.18) {$M_\alpha$};
        \end{scope}
    \end{tikzpicture}%
}

% Inverse-map contraction for an injective simple tensor: contracting the
% inverse tensor with the tensor through the physical index identifies the
% two pairs of virtual legs (the matrix units on the bond space).
\newcommand{\TNMPDOInverseContraction}{%
    \begin{tikzpicture}[tn picture]
        \TNTensor{invT}{(0,0.55)}
        \TNTensor{invB}{(0,-0.55)}
        \TN@vleg{invT}{west}{-0.60,0}
        \TN@vleg{invT}{east}{0.60,0}
        \TN@vleg{invB}{west}{-0.60,0}
        \TN@vleg{invB}{east}{0.60,0}
        \TN@ppath{ (invT.south) -- (invB.north);}
        \node at (-0.50,0.88) {$K^{-1}$};
        \node at (-0.42,-0.88) {$K$};
        \node at (1.35,0) {$=$};
        \TN@vpath{ (1.90,0.55) -- (2.15,0.55)
        .. controls (2.70,0.55) and (2.70,-0.55) .. (2.15,-0.55)
        -- (1.90,-0.55);}
        \TN@vpath{ (3.90,0.55) -- (3.65,0.55)
        .. controls (3.10,0.55) and (3.10,-0.55) .. (3.65,-0.55)
        -- (3.90,-0.55);}
    \end{tikzpicture}%
}

% The normalized three-site marginal of a four-site periodic MPO.  The first
% three sites retain their two physical legs, while tracing the fourth site is
% absorbed into the virtual closing operator R_4.
\newcommand{\TNMPDONormalizedFourSiteTail}{%
    \begin{tikzpicture}[tn picture]
        \node at (-3.72,0) {$\sigma_3^{(4)}({\cal K})$};
        \node at (-2.76,0) {$=$};

        \TNTensor{tnTailOne}{(-1.48,0)}
        \TNTensor{tnTailTwo}{(0,0)}
        \TNTensor{tnTailThree}{(1.48,0)}
        \TN@vpath{ (tnTailOne.east) -- (tnTailTwo.west);}
        \TN@vpath{ (tnTailTwo.east) -- (tnTailThree.west);}

        \foreach \n in {tnTailOne,tnTailTwo,tnTailThree} {
            \TN@pleg{\n}{north}{0,0.48}
            \TN@pleg{\n}{south}{0,-0.48}
        }
        \TN@uplabel{tnTailOne}{0.72}{i_1}
        \TN@uplabel{tnTailTwo}{0.72}{i_2}
        \TN@uplabel{tnTailThree}{0.72}{i_3}
        \TN@downlabel{tnTailOne}{0.72}{j_1}
        \TN@downlabel{tnTailTwo}{0.72}{j_2}
        \TN@downlabel{tnTailThree}{0.72}{j_3}
        \node[anchor=east] at (-1.62,0.25) {${\cal K}$};
        \node[anchor=east] at (-0.14,0.25) {${\cal K}$};
        \node[anchor=east] at (1.34,0.25) {${\cal K}$};

        \TNComponent{tnFourSiteTail}{(0,-1.42)}{R_4}
        \TN@vtracepath{ (tnTailOne.west) -- ++(-0.48,0)
            |- ($(tnFourSiteTail.east)+(0,-0.40)$) -- (tnFourSiteTail.east);}
        \TN@vtracepath{ (tnTailThree.east) -- ++(0.48,0)
            |- ($(tnFourSiteTail.west)+(0,0.38)$) -- (tnFourSiteTail.west);}

    \end{tikzpicture}%
}

% Sectorwise comparison between the inverse-map contraction and the Hayashi
% product in the basis adapted to the middle-site direct sum.
\newcommand{\TNMPDOHayashiSectorComparison}{%
    \begin{tikzpicture}[tn picture]
        \node at (-4.55,0) {$R_{\beta_3,\alpha_1}$};
        \node at (-3.65,0) {$\cdot$};

        \TNTensor{tnSectorK}{(-2.45,0)}
        \TN@vleg{tnSectorK}{west}{-0.72,0}
        \TN@vleg{tnSectorK}{east}{0.72,0}
        \node[anchor=east] at (-3.22,0) {$\beta_1$};
        \node[anchor=west] at (-1.68,0) {$\alpha_3$};
        \node[anchor=east] at (-2.58,0.25) {${\cal K}$};

        \TNMap{tnSectorU}{(-2.45,0.88)}{U_B}
        \TNMap{tnSectorUd}{(-2.45,-0.88)}{U_B^\dagger}
        \TN@ppath{ (tnSectorK.north) -- (tnSectorU.south);}
        \TN@ppath{ (tnSectorU.north) -- ++(0,0.34);}
        \TN@ppath{ (tnSectorK.south) -- (tnSectorUd.north);}
        \TN@ppath{ (tnSectorUd.south) -- ++(0,-0.34);}
        \node[anchor=west] at (-1.92,0.88) {$k;l,r$};
        \node[anchor=west] at (-1.92,-0.88) {$k;l',r'$};

        \node at (-0.60,0) {$=$};
        \node at (0.10,0) {$p_k$};
        \TN@factor[minimum width=1.00cm, minimum height=0.62cm]
            {tnSectorA}{(1.25,0)}{A^{(k)}_{\alpha_1,\beta_1}}
        \node at (2.15,0) {$\otimes$};
        \TN@factor[minimum width=1.00cm, minimum height=0.62cm]
            {tnSectorB}{(3.25,0)}{B^{(k)}_{\alpha_3,\beta_3}}
        \node at (1.25,0.62) {$l,l'$};
        \node at (3.25,0.62) {$r,r'$};
    \end{tikzpicture}%
}

% Sector factorization: in the middle-site basis selected by U_B, every physical
% slice of the injective tensor is a direct sum of the sector products
% l_k \otimes r_k.  The placement of U_B and U_B^\dagger follows the theorem
% orientation U_B K U_B^\dagger.
\newcommand{\TNMPDOSectorFactorization}{%
    \begin{tikzpicture}[tn picture]
        \TNTensor{sfK}{(-2.65,0)}
        \TN@vleg{sfK}{west}{-0.72,0}
        \TN@vleg{sfK}{east}{0.72,0}
        \node[anchor=east] at (-3.42,0) {$\beta_1$};
        \node[anchor=west] at (-1.88,0) {$\alpha_3$};
        \node[anchor=east] at (-2.78,0.25) {${\cal K}$};

        \TNMap{sfU}{(-2.65,0.88)}{U_B}
        \TNMap{sfUd}{(-2.65,-0.88)}{U_B^\dagger}
        \TN@ppath{ (sfK.north) -- (sfU.south);}
        \TN@ppath{ (sfU.north) -- ++(0,0.34);}
        \TN@ppath{ (sfK.south) -- (sfUd.north);}
        \TN@ppath{ (sfUd.south) -- ++(0,-0.34);}

        \node at (-0.95,0) {$=$};
        \node at (-0.10,0) {$\displaystyle\bigoplus_k$};
        \TNTensor{sfL}{(1.10,0)}
        \TNTensor{sfR}{(2.75,0)}
        \TN@vpath{ (sfL.west) -- ++(-0.62,0);}
        \node at (1.93,0) {$\otimes$};
        \TN@vpath{ (sfR.east) -- ++(0.62,0);}
        \TN@pleg{sfL}{north}{0,0.50}
        \TN@pleg{sfL}{south}{0,-0.50}
        \TN@pleg{sfR}{north}{0,0.50}
        \TN@pleg{sfR}{south}{0,-0.50}
        \node at (1.10,0.78) {$(l_k)_{\beta_1}$};
        \node at (2.75,0.78) {$(r_k)_{\alpha_3}$};
    \end{tikzpicture}%
}

% Local-orthogonality contraction: sandwiching the ket and bra legs of one
% tensor between the orthogonal projectors P_1 and P_2, in either order,
% gives zero, as does the neighboring two-site contraction with P_1 dressing
% both physical legs of the first site and P_2 both legs of the second.
\newcommand{\TNMPDOLocalOrthogonality}{%
    \begin{tikzpicture}[tn picture]
        % One site, P_1 on the ket leg and P_2 on the bra leg.
        \TNTensor{loA}{(0,0)}
        \TN@vleg{loA}{west}{-0.50,0}
        \TN@vleg{loA}{east}{0.50,0}
        \TN@pleg{loA}{north}{0,1.00}
        \TN@pleg{loA}{south}{0,-1.00}
        \TN@insertion[label=left:$P_1$]{loAu}{(0,0.55)}
        \TN@insertion[label=left:$P_2$]{loAd}{(0,-0.55)}
        \node[anchor=west] at (0.12,-0.30) {${\cal K}$};
        \node at (1.10,0) {$=$};
        % One site, P_2 on the ket leg and P_1 on the bra leg.
        \TNTensor{loB}{(2.20,0)}
        \TN@vleg{loB}{west}{-0.50,0}
        \TN@vleg{loB}{east}{0.50,0}
        \TN@pleg{loB}{north}{0,1.00}
        \TN@pleg{loB}{south}{0,-1.00}
        \TN@insertion[label=left:$P_2$]{loBu}{(2.20,0.55)}
        \TN@insertion[label=left:$P_1$]{loBd}{(2.20,-0.55)}
        \node[anchor=west] at (2.32,-0.30) {${\cal K}$};
        \node at (3.30,0) {$=$};
        % Two neighboring sites, P_1 dressing both physical legs of the
        % first site and P_2 both physical legs of the second.
        \TNTensor{loCa}{(4.40,0)}
        \TNTensor{loCb}{(5.60,0)}
        \TN@vpath{ (loCa.east) -- (loCb.west);}
        \TN@vleg{loCa}{west}{-0.50,0}
        \TN@vleg{loCb}{east}{0.50,0}
        \TN@pleg{loCa}{north}{0,1.00}
        \TN@pleg{loCa}{south}{0,-1.00}
        \TN@pleg{loCb}{north}{0,1.00}
        \TN@pleg{loCb}{south}{0,-1.00}
        \TN@insertion[label=left:$P_1$]{loCau}{(4.40,0.55)}
        \TN@insertion[label=left:$P_1$]{loCad}{(4.40,-0.55)}
        \TN@insertion[label=right:$P_2$]{loCbu}{(5.60,0.55)}
        \TN@insertion[label=right:$P_2$]{loCbd}{(5.60,-0.55)}
        \node[anchor=west] at (4.52,-0.30) {${\cal K}$};
        \node[anchor=west] at (5.72,-0.30) {${\cal K}$};
        \node at (6.65,0) {$=\;0$};
    \end{tikzpicture}%
}

% Inverse-tensor contraction for an injective MPS tensor: contracting the
% inverse tensor with the tensor through the shared physical index identifies
% the virtual legs of the inverse with the corresponding virtual legs of the
% tensor.
\newcommand{\TNMPSInverseContraction}{%
    \begin{tikzpicture}[tn picture]
        \TNTensor{mpsiT}{(0,0.55)}
        \TNTensor{mpsiB}{(0,-0.55)}
        \TN@vleg{mpsiT}{west}{-0.60,0}
        \TN@vleg{mpsiT}{east}{0.60,0}
        \TN@vleg{mpsiB}{west}{-0.60,0}
        \TN@vleg{mpsiB}{east}{0.60,0}
        \TN@ppath{ (mpsiT.south) -- (mpsiB.north);}
        \node at (-0.52,0.88) {$A^{-1}$};
        \node at (-0.42,-0.88) {$A$};
        \node at (1.35,0) {$=$};
        \TN@vpath{ (1.90,0.55) -- (2.15,0.55)
        .. controls (2.70,0.55) and (2.70,-0.55) .. (2.15,-0.55)
        -- (1.90,-0.55);}
        \TN@vpath{ (3.90,0.55) -- (3.65,0.55)
        .. controls (3.10,0.55) and (3.10,-0.55) .. (3.65,-0.55)
        -- (3.90,-0.55);}
    \end{tikzpicture}%
}

% Coefficient form of a basis of normal tensors: the gauge X, one basis
% tensor with its physical leg, and X^{-1} on the virtual chain, threaded by
% the block-index lines j and q; the coefficient tensor mu hangs on the two
% index lines, so the direct sums are drawn as index lines.
\newcommand{\TNBNTDecomposition}{%
    \begin{tikzpicture}[tn picture]
        \TN@insertion[label=above:$X$]{bntX}{(0,0)}
        \TNTensor{bntA}{(1.50,0)}
        \TN@insertion[label=above:$X^{-1}$]{bntXi}{(3.00,0)}
        \TN@vleg{bntX}{west}{-0.55,0}
        \TN@vconnect{bntX}{east}{bntA}{west}
        \TN@vconnect{bntA}{east}{bntXi}{west}
        \TN@vleg{bntXi}{east}{0.55,0}
        \TN@pleg{bntA}{north}{0,0.46}
        \TN@label{bntA.north}{(0,0.70)}{i}
        \node[anchor=west] at (1.62,-0.26) {$A_j$};
        \node[anchor=east] at (-0.70,0) {$\alpha$};
        \node[anchor=west] at (3.70,0) {$\beta$};
        \TN@vpath{(-0.60,-0.62) -- (3.60,-0.62)}
        \TN@vpath{(-0.60,-1.02) -- (3.60,-1.02)}
        \node[anchor=east] at (-0.70,-0.62) {$j$};
        \node[anchor=east] at (-0.70,-1.02) {$q$};
        \TN@vpath{(bntX.south) -- (0,-1.02)}
        \TN@vpath{(bntA.south) -- (1.50,-0.62)}
        \TN@vpath{(bntXi.south) -- (3.00,-1.02)}
        \TNComponent{bntMu}{(0.75,-1.60)}{\mu}
        \TN@vnorthport{bntMuNLeft}{bntMu}{0.15}
        \TN@vnorthport{bntMuNRight}{bntMu}{0.85}
        \TN@vpath{(bntMuNLeft) -- (0.63,-0.62)}
        \TN@vpath{(bntMuNRight) -- (0.87,-1.02)}
        \TN@anonymousjunction{(0,-0.62)}
        \TN@anonymousjunction{(0,-1.02)}
        \TN@anonymousjunction{(1.50,-0.62)}
        \TN@anonymousjunction{(3.00,-0.62)}
        \TN@anonymousjunction{(3.00,-1.02)}
        \TN@anonymousjunction{(0.63,-0.62)}
        \TN@anonymousjunction{(0.87,-1.02)}
    \end{tikzpicture}%
}

% Zero-correlation-length idempotence for the physical-trace transfer: the
% trace-closed tensor equals the virtual product of two trace-closed copies.
\newcommand{\TNMPDOZCLIdempotence}{%
    \begin{TNEquationRow}[compact]
        \TNTerm{%
            \TNMPOSite{zciM}{(0,0)}{M}
            \TNOpenVirtualWest{zciMWest}
            \TNOpenVirtualEast{zciMEast}
            \TNTracePhysicalRight{zciMKet}{zciMBra}
            \node[below] at (zciM) {$\mathcal T_M$};}
        \TNRelation{=}
        \TNTerm{%
            \TNMPOSite{zciA}{(0,0)}{M}
            \TNMPOSite{zciB}{(1.25,0)}{M}
            \TNOpenVirtualWest{zciAWest}
            \TNConnectVirtual{zciAEast}{zciBWest}
            \TNOpenVirtualEast{zciBEast}
            \TNTracePhysicalRight{zciAKet}{zciABra}
            \TNTracePhysicalRight{zciBKet}{zciBBra}
            \node[below] at (zciA) {$\mathcal T_M$};
            \node[below] at (zciB) {$\mathcal T_M$};}
    \end{TNEquationRow}%
}

% The two fixed-final-label fusion trees.  The dashed bridge records the
% F-change of multiplicity basis that remains to be constructed.
\newcommand{\TNMPDOFixedFinalFusionBracketings}{%
    \begin{TNDiagram}[normal]
        \path[use as bounding box] (-4.70,-2.05) rectangle (4.70,1.55);
        \TNCofusionMap{ffLeftLower}{(-3.35,-0.72)}{U_{\alpha,\beta}^{\delta}}
        \TNCofusionMap{ffLeftUpper}{(-2.80,0.55)}{U_{\delta,\gamma}^{\varepsilon}}
        \TNConnectVirtual{ffLeftLowerCombined}{ffLeftUpperFactorOne}
        \TNOpenVirtualSouth{ffLeftLowerFactorOne}
        \TNOpenVirtualSouth{ffLeftLowerFactorTwo}
        \TNOpenVirtualSouth{ffLeftUpperFactorTwo}
        \TNOpenVirtualNorth{ffLeftUpperCombined}
        \node[below] at (ffLeftLowerFactorOneOpen) {$\alpha$};
        \node[below] at (ffLeftLowerFactorTwoOpen) {$\beta$};
        \node[below] at (ffLeftUpperFactorTwoOpen) {$\gamma$};
        \node[above] at (ffLeftUpperCombinedOpen) {$\varepsilon$};
        \node at (-3.00,-1.72) {$((\alpha\beta)\gamma)$};

        \TNCofusionMap{ffRightLower}{(3.35,-0.72)}{U_{\beta,\gamma}^{\eta}}
        \TNCofusionMap{ffRightUpper}{(2.80,0.55)}{U_{\alpha,\eta}^{\varepsilon}}
        \TNConnectVirtual{ffRightLowerCombined}{ffRightUpperFactorTwo}
        \TNOpenVirtualSouth{ffRightUpperFactorOne}
        \TNOpenVirtualSouth{ffRightLowerFactorOne}
        \TNOpenVirtualSouth{ffRightLowerFactorTwo}
        \TNOpenVirtualNorth{ffRightUpperCombined}
        \node[below] at (ffRightUpperFactorOneOpen) {$\alpha$};
        \node[below] at (ffRightLowerFactorOneOpen) {$\beta$};
        \node[below] at (ffRightLowerFactorTwoOpen) {$\gamma$};
        \node[above] at (ffRightUpperCombinedOpen) {$\varepsilon$};
        \node at (3.00,-1.72) {$(\alpha(\beta\gamma))$};

        \TNComparisonArrow{ffBracketingChange}{(ffLeftUpper.east)}
            {(ffRightUpper.west)}{F_{\varepsilon}^{\alpha\beta\gamma}}
    \end{TNDiagram}%
}

% The source-oriented F-move.  The left tree is expanded in the basis of
% right trees.  The multiplicity labels and the direction of F agree with
% equation (Fmove) of arXiv:1511.08090.
\newcommand{\TNMPDOPrintedFMove}{%
    \begin{TNDiagram}[normal]
            \TNCofusionMap{pfLeftLower}{(-4.25,-0.55)}{X^{e}_{ab,\mu}}
            \TNCofusionMap{pfLeftUpper}{(-3.70,0.70)}{X^{d}_{ec,\nu}}
            \TNConnectVirtual{pfLeftLowerCombined}{pfLeftUpperFactorOne}
            \TNOpenVirtualSouth{pfLeftLowerFactorOne}
            \TNOpenVirtualSouth{pfLeftLowerFactorTwo}
            \TNOpenVirtualSouth{pfLeftUpperFactorTwo}
            \TNOpenVirtualNorth{pfLeftUpperCombined}
            \node[below] at (pfLeftLowerFactorOneOpen) {$a$};
            \node[below] at (pfLeftLowerFactorTwoOpen) {$b$};
            \node[below] at (pfLeftUpperFactorTwoOpen) {$c$};
            \node[above] at (pfLeftUpperCombinedOpen) {$d$};
        \node at (-1.55,0) {$=\displaystyle\sum_{f,\lambda,\sigma}$};
        \node at (0.25,0)
            {$\displaystyle(F^{abc}_{d})^{f\lambda\sigma}_{e\mu\nu}$};
            \TNCofusionMap{pfRightLower}{(3.15,-0.55)}{X^{f}_{bc,\lambda}}
            \TNCofusionMap{pfRightUpper}{(2.60,0.70)}{X^{d}_{af,\sigma}}
            \TNConnectVirtual{pfRightLowerCombined}{pfRightUpperFactorTwo}
            \TNOpenVirtualSouth{pfRightUpperFactorOne}
            \TNOpenVirtualSouth{pfRightLowerFactorOne}
            \TNOpenVirtualSouth{pfRightLowerFactorTwo}
            \TNOpenVirtualNorth{pfRightUpperCombined}
            \node[below] at (pfRightUpperFactorOneOpen) {$a$};
            \node[below] at (pfRightLowerFactorOneOpen) {$b$};
            \node[below] at (pfRightLowerFactorTwoOpen) {$c$};
            \node[above] at (pfRightUpperCombinedOpen) {$d$};
    \end{TNDiagram}%
}

% The fixed-final comparison argument in three graphical steps.  A
% positive-length simultaneous word selector isolates one final label, the
% resulting comparison is block diagonal in that label, and injectivity
% identifies each diagonal corner with F tensored with the identity on the
% final bond space.  The last line records two-sided unitarity, but deliberately
% contains no four-fold reassociation or pentagon claim.
\newcommand{\TNMPDOFixedFinalComparisonUnitary}{%
    \begin{TNDiagram}[normal]
        \path[use as bounding box] (-5.90,-3.45) rectangle (5.90,3.25);
        \node[anchor=east] at (-5.10,2.35) {\textup{selector}, $S>0$};
        \TNMPOSite{ffSelL}{(-3.90,2.35)}{}
        \TNMPOSite{ffSelM}{(-2.55,2.35)}{M_{\varepsilon'}^{w}}
        \TNMPOSite{ffSelR}{(-1.20,2.35)}{}
        \TNConnectVirtual{ffSelLEast}{ffSelMWest}
        \TNConnectVirtual{ffSelMEast}{ffSelRWest}
        \TNOpenVirtualWest{ffSelLWest}
        \TNOpenVirtualEast{ffSelREast}
        \node[above] at (-2.55,2.78)
            {$\sum_{|w|=S}c_{\varepsilon}(w)(\,\cdot\,)^{w}$};
        \TNMorphismArrow{ffSelectorArrow}{(-0.35,2.35)}{(0.60,2.35)}{}
        \node[anchor=west] at (0.78,2.35)
            {$\delta_{\varepsilon,\varepsilon'}1_{D_\varepsilon}$};

        \node[anchor=east] at (-5.10,0.30) {\textup{sector cut}};
        \TNBinaryFusionTree{ffCutLeft}{(-2.65,0.20)}
            {\alpha}{\beta}{\gamma}{U^{\mathrm L}}
        \TNBinaryFusionTree{ffCutRight}{(2.65,0.20)}
            {\alpha}{\beta}{\gamma}{U^{\mathrm R}}
        \TNOpenVirtualSouth{ffCutLeftLowerCombined}
        \TNOpenVirtualSouth{ffCutRightLowerCombined}
        \TNComparisonArrow{ffSectorComparison}{(ffCutLeftUpper.east)}
            {(ffCutRightUpper.west)}
            {C^{\varepsilon,\varepsilon'}_{\alpha,\beta,\gamma}}
        \node[anchor=west] at (4.25,0.20)
            {$=0\quad(\varepsilon\ne\varepsilon')$};

        \node[anchor=east] at (-5.10,-2.15) {\textup{diagonal corner}};
        \TNExpression{ffLeftSpace}{(-3.35,-2.15)}
            {\mathcal H^{\mathrm L}_{\varepsilon}\otimes\C^{D_\varepsilon}}
        \TNMap{ffF}{(0,-2.15)}{F_\varepsilon\otimes1_{D_\varepsilon}}
        \TNExpression{ffRightSpace}{(3.35,-2.15)}
            {\mathcal H^{\mathrm R}_{\varepsilon}\otimes\C^{D_\varepsilon}}
        \TNMorphismArrow{ffDiagonalLeft}{(ffLeftSpace.east)}{(ffF.west)}{}
        \TNMorphismArrow{ffDiagonalRight}{(ffF.east)}{(ffRightSpace.west)}{}
        \node at (0,-3.05)
            {$F_\varepsilon F_\varepsilon^\dagger=1,\qquad
              F_\varepsilon^\dagger F_\varepsilon=1$};
    \end{TNDiagram}%
}

% The associahedron for four ordinary bond-coordinate factors.  The upper
% three-edge route and lower two-edge route have the same endpoints.  This is
% only Cartesian-product reassociation; the edges do not represent F-matrices.
\newcommand{\TNMPDOFourfoldBondReassociationPentagon}{%
    \begin{TNDiagram}[normal]
        \TNExpression{assocL}{(-5.25,0)}{(((A\times B)\times C)\times D)}
        \TNExpression{assocLI}{(-3.10,2.05)}{((A\times(B\times C))\times D)}
        \TNExpression{assocM}{(1.20,2.05)}{(A\times((B\times C)\times D))}
        \TNExpression{assocR}{(5.25,0)}{(A\times(B\times(C\times D)))}
        \TNExpression{assocP}{(0,-2.05)}{((A\times B)\times(C\times D))}
        \TNMorphismArrow{assocRone}{(assocL)}{(assocLI)}{r_1}
        \TNMorphismArrow{assocRtwo}{(assocLI)}{(assocM)}{r_2}
        \TNMorphismArrow{assocRthree}{(assocM)}{(assocR)}{r_3}
        \TNMorphismArrow{assocSone}{(assocL)}{(assocP)}{s_1}
        \TNMorphismArrow{assocStwo}{(assocP)}{(assocR)}{s_2}
    \end{TNDiagram}%
}

% The five genuine fusion trees in the MPO pentagon.  Unlike the Cartesian
% reassociation diagram above, every edge is a printed-orientation F-move and
% every internal label is a fusion-sector label.  The arrangement follows
% Figure `pentagon` (n13.pdf) of arXiv:1511.08090.
\newcommand{\TNMPDOCompleteZipperFusionPentagon}{%
    \begin{TNDiagram}[normal]
        \path[use as bounding box] (-6.25,-4.30) rectangle (6.25,4.10);
        % (((a b) c) d)
        \TNFusionMap{pLf}{(-5.42,0.22)}{}
        \TNFusionMap{pLg}{(-5.00,-0.55)}{}
        \TNFusionMap{pLe}{(-4.56,-1.32)}{}
        \TNOpenVirtualNorth{pLfFactorOne}
        \TNOpenVirtualNorth{pLfFactorTwo}
        \TNConnectVirtual{pLfCombined}{pLgFactorOne}
        \TNOpenVirtualNorth{pLgFactorTwo}
        \TNConnectVirtual{pLgCombined}{pLeFactorOne}
        \TNOpenVirtualNorth{pLeFactorTwo}
        \TNOpenVirtualSouth{pLeCombined}
        \foreach \n/\t in {pLfFactorOneOpen/a,pLfFactorTwoOpen/b,
            pLgFactorTwoOpen/c,pLeFactorTwoOpen/d}
            \node[above] at (\n) {$\t$};
        \node[left] at ($(pLfOut)!0.5!(pLgLeft)$) {$f$};
        \node[left] at ($(pLgOut)!0.5!(pLeLeft)$) {$g$};
        \node[below] at (pLeCombinedOpen) {$e$};

        % ((a (b c)) d)
        \TNFusionMap{pLIh}{(-3.38,2.99)}{}
        \TNFusionMap{pLIg}{(-3.78,2.22)}{}
        \TNFusionMap{pLIe}{(-3.28,1.45)}{}
        \TNOpenVirtualNorth{pLIhFactorOne}
        \TNOpenVirtualNorth{pLIhFactorTwo}
        \TNOpenVirtualNorth{pLIgFactorOne}
        \TNConnectVirtual{pLIhCombined}{pLIgFactorTwo}
        \TNConnectVirtual{pLIgCombined}{pLIeFactorOne}
        \TNOpenVirtualNorth{pLIeFactorTwo}
        \TNOpenVirtualSouth{pLIeCombined}
        \foreach \n/\t in {pLIgFactorOneOpen/a,pLIhFactorOneOpen/b,
            pLIhFactorTwoOpen/c,pLIeFactorTwoOpen/d}
            \node[above] at (\n) {$\t$};
        \node[right] at ($(pLIhOut)!0.5!(pLIgRight)$) {$h$};
        \node[left] at ($(pLIgOut)!0.5!(pLIeLeft)$) {$g$};
        \node[below] at (pLIeCombinedOpen) {$e$};

        % a ((b c) d)
        \TNFusionMap{pMh}{(0.92,2.99)}{}
        \TNFusionMap{pMi}{(1.34,2.22)}{}
        \TNFusionMap{pMe}{(0.84,1.45)}{}
        \TNOpenVirtualNorth{pMhFactorOne}
        \TNOpenVirtualNorth{pMhFactorTwo}
        \TNConnectVirtual{pMhCombined}{pMiFactorOne}
        \TNOpenVirtualNorth{pMiFactorTwo}
        \TNOpenVirtualNorth{pMeFactorOne}
        \TNConnectVirtual{pMiCombined}{pMeFactorTwo}
        \TNOpenVirtualSouth{pMeCombined}
        \foreach \n/\t in {pMeFactorOneOpen/a,pMhFactorOneOpen/b,
            pMhFactorTwoOpen/c,pMiFactorTwoOpen/d}
            \node[above] at (\n) {$\t$};
        \node[left] at ($(pMhOut)!0.5!(pMiLeft)$) {$h$};
        \node[right] at ($(pMiOut)!0.5!(pMeRight)$) {$i$};
        \node[below] at (pMeCombinedOpen) {$e$};

        % a (b (c d))
        \TNFusionMap{pRj}{(5.42,0.22)}{}
        \TNFusionMap{pRi}{(5.00,-0.55)}{}
        \TNFusionMap{pRe}{(4.56,-1.32)}{}
        \TNOpenVirtualNorth{pRjFactorOne}
        \TNOpenVirtualNorth{pRjFactorTwo}
        \TNOpenVirtualNorth{pRiFactorOne}
        \TNConnectVirtual{pRjCombined}{pRiFactorTwo}
        \TNOpenVirtualNorth{pReFactorOne}
        \TNConnectVirtual{pRiCombined}{pReFactorTwo}
        \TNOpenVirtualSouth{pReCombined}
        \foreach \n/\t in {pReFactorOneOpen/a,pRiFactorOneOpen/b,
            pRjFactorOneOpen/c,pRjFactorTwoOpen/d}
            \node[above] at (\n) {$\t$};
        \node[right] at ($(pRjOut)!0.5!(pRiRight)$) {$j$};
        \node[right] at ($(pRiOut)!0.5!(pReRight)$) {$i$};
        \node[below] at (pReCombinedOpen) {$e$};

        % (a b) (c d)
        \TNFusionMap{pPf}{(-0.72,-2.63)}{}
        \TNFusionMap{pPj}{(0.72,-2.63)}{}
        \TNFusionMap{pPe}{(0,-3.48)}{}
        \TNOpenVirtualNorth{pPfFactorOne}
        \TNOpenVirtualNorth{pPfFactorTwo}
        \TNOpenVirtualNorth{pPjFactorOne}
        \TNOpenVirtualNorth{pPjFactorTwo}
        \TNConnectVirtual{pPfCombined}{pPeFactorOne}
        \TNConnectVirtual{pPjCombined}{pPeFactorTwo}
        \TNOpenVirtualSouth{pPeCombined}
        \foreach \n/\t in {pPfFactorOneOpen/a,pPfFactorTwoOpen/b,
            pPjFactorOneOpen/c,pPjFactorTwoOpen/d}
            \node[above] at (\n) {$\t$};
        \node[left] at ($(pPfOut)!0.5!(pPeLeft)$) {$f$};
        \node[right] at ($(pPjOut)!0.5!(pPeRight)$) {$j$};
        \node[below] at (pPeCombinedOpen) {$e$};

        \TNMorphismArrow{pentagonOne}{(-4.48,1.10)}{(-3.95,1.73)}{F^{abc}_{g}}
        \TNMorphismArrow{pentagonTwo}{(-2.20,2.72)}{(-0.25,2.72)}{F^{ahd}_{e}}
        \TNMorphismArrow{pentagonThree}{(1.80,1.78)}{(4.15,1.08)}{F^{bcd}_{i}}
        \TNMorphismArrow{pentagonFour}{(-3.95,-1.05)}{(-1.02,-2.15)}{F^{fcd}_{e}}
        \TNMorphismArrow{pentagonFive}{(1.02,-2.15)}{(3.95,-1.05)}{F^{abj}_{e}}
    \end{TNDiagram}%
}

% Local BNT fusion identity from CPSV16, eq. (Ualphabeta).  The left-hand
% isometry combines the two incoming bond factors, while its adjoint separates
% the outgoing factors.  The right-hand side records the weighted final-sector
% direct sum without introducing an F-move.
\newcommand{\TNMPDOBNTFusionIdentity}{%
    \begin{TNDiagram}[compact]
            \TNSplitMap{bfiU}{(-3.65,0)}{U_{\alpha,\beta}}
            \TNStackedMPOProduct{bfiProduct}{(-2.35,0)}{M_\alpha}{M_\beta}{}
            \TNMergeMap{bfiUd}{(-1.05,0)}{U_{\alpha,\beta}^{\dagger}}
            \TNConnectVirtual{bfiUFactorOne}{bfiProductUpperWest}
            \TNConnectVirtual{bfiUFactorTwo}{bfiProductLowerWest}
            \TNConnectVirtual{bfiProductUpperEast}{bfiUdFactorOne}
            \TNConnectVirtual{bfiProductLowerEast}{bfiUdFactorTwo}
            \TNOpenVirtualWest{bfiUCombined}
            \TNOpenVirtualEast{bfiUdCombined}
            \TNOpenPhysicalNorth{bfiProductUpperKet}
            \TNOpenPhysicalSouth{bfiProductLowerBra}
            \node[above] at (bfiProductUpperKetOpen) {$i$};
            \node[below] at (bfiProductLowerBraOpen) {$j$};
        \node at (0.45,0) {$=\displaystyle\bigoplus_\gamma$};
            \TNWeightedMPOBlock{bfiBlock}{(2.30,0)}
                {M_\gamma}{M_\gamma}{}{\chi_{\alpha,\beta,\gamma}}
            \TNOpenVirtualWest{bfiBlockWeightW}
            \TNOpenVirtualEast{bfiBlockWeightedEast}
            \TNOpenVirtualWest{bfiBlockLowerWest}
            \TNOpenVirtualEast{bfiBlockLowerEast}
            \TNOpenPhysicalNorth{bfiBlockUpperKet}
            \TNOpenPhysicalSouth{bfiBlockLowerBra}
    \end{TNDiagram}%
}

% The two unweighted zipper directions and the resulting reconstruction of a
% product letter.  This follows the graphical arrangement of the zipper
% figures in Bultinck et al., arXiv:1511.08090, while using the notation of
% the present MPDO fusion family.  The final row uses U^\dagger U = 1 only;
% it deliberately does not depict U U^\dagger = 1.
\newcommand{\TNMPDOUnweightedZipperReconstruction}{%
    \begin{TNEquationRow}[compact]
      \TNTerm{\TNSplitMap{uzOneU}{(-3.80,2.20)}{U_{\alpha,\beta}}
            \TNStackedMPOProduct{uzOneP}{(-2.65,2.20)}{M_\alpha}{M_\beta}{}
            \TNConnectVirtual{uzOneUFactorOne}{uzOnePUpperWest}
            \TNConnectVirtual{uzOneUFactorTwo}{uzOnePLowerWest}
            \TNOpenVirtualWest{uzOneUCombined}
            \TNOpenVirtualEast{uzOnePUpperEast}
            \TNOpenVirtualEast{uzOnePLowerEast}
        \node at (-1.55,2.20) {$=$};
        \TNMPOSite{uzOneH}{(-0.35,2.20)}{H_{\alpha,\beta}^{ij}}
            \TNSplitMap{uzOneR}{(1.15,2.20)}{U_{\alpha,\beta}}
            \TNConnectVirtual{uzOneHEast}{uzOneRCombined}
            \TNOpenVirtualWest{uzOneHWest}
            \TNOpenVirtualEast{uzOneRFactorOne}
            \TNOpenVirtualEast{uzOneRFactorTwo}

        \TNStackedMPOProduct{uzTwoP}{(-3.80,0)}{M_\alpha}{M_\beta}{}
            \TNMergeMap{uzTwoU}{(-2.65,0)}{U_{\alpha,\beta}^{\dagger}}
            \TNConnectVirtual{uzTwoPUpperEast}{uzTwoUFactorOne}
            \TNConnectVirtual{uzTwoPLowerEast}{uzTwoUFactorTwo}
            \TNOpenVirtualWest{uzTwoPUpperWest}
            \TNOpenVirtualWest{uzTwoPLowerWest}
            \TNOpenVirtualEast{uzTwoUCombined}
        \node at (-1.55,0) {$=$};
        \TNMergeMap{uzTwoL}{(-0.35,0)}{U_{\alpha,\beta}^{\dagger}}
            \TNMPOSite{uzTwoH}{(1.15,0)}{H_{\alpha,\beta}^{ij}}
            \TNConnectVirtual{uzTwoLCombined}{uzTwoHWest}
            \TNOpenVirtualWest{uzTwoLFactorOne}
            \TNOpenVirtualWest{uzTwoLFactorTwo}
            \TNOpenVirtualEast{uzTwoHEast}

        \TNStackedMPOProduct{uzThreeP}{(-3.80,-2.20)}{M_\alpha}{M_\beta}{}
            \TNOpenVirtualWest{uzThreePUpperWest}
            \TNOpenVirtualWest{uzThreePLowerWest}
            \TNOpenVirtualEast{uzThreePUpperEast}
            \TNOpenVirtualEast{uzThreePLowerEast}
        \node at (-2.65,-2.20) {$=$};
        \TNMergeMap{uzThreeL}{(-1.45,-2.20)}{U_{\alpha,\beta}^{\dagger}}
            \TNMPOSite{uzThreeH}{(0,-2.20)}{H_{\alpha,\beta}^{ij}}
            \TNSplitMap{uzThreeR}{(1.45,-2.20)}{U_{\alpha,\beta}}
            \TNConnectVirtual{uzThreeLCombined}{uzThreeHWest}
            \TNConnectVirtual{uzThreeHEast}{uzThreeRCombined}
            \TNOpenVirtualWest{uzThreeLFactorOne}
            \TNOpenVirtualWest{uzThreeLFactorTwo}
            \TNOpenVirtualEast{uzThreeRFactorOne}
            \TNOpenVirtualEast{uzThreeRFactorTwo}}
    \end{TNEquationRow}%
}

% Closing an L-site chain of the local BNT fusion identity.  On the right the
% multiplicity loop and the M_gamma loop are disjoint, so their traces factor.
\newcommand{\TNMPDOFusionTracePower}{%
    \begin{TNDiagram}[compact]
            \TNStackedMPOProduct{ftpLeftA}{(-5.20,0)}{M_\alpha}{M_\beta}{}
            \TNStackedMPOProduct{ftpLeftB}{(-4.00,0)}{M_\alpha}{M_\beta}{}
            \TNStackedMPOProduct{ftpLeftC}{(-2.80,0)}{M_\alpha}{M_\beta}{}
            \TNConnectVirtual{ftpLeftAUpperEast}{ftpLeftBUpperWest}
            \TNConnectVirtual{ftpLeftBUpperEast}{ftpLeftCUpperWest}
            \TNConnectVirtual{ftpLeftALowerEast}{ftpLeftBLowerWest}
            \TNConnectVirtual{ftpLeftBLowerEast}{ftpLeftCLowerWest}
            \TNTraceVirtualAbove{ftpLeftAUpperWest}{ftpLeftCUpperEast}
            \TNTraceVirtualBelow{ftpLeftALowerWest}{ftpLeftCLowerEast}
            \node at (-4.00,-1.45) {$L\text{ sites}$};
        \node at (-1.05,0) {$=\displaystyle\sum_\gamma$};
            \TNStackedMPOProduct{ftpRightA}{(0.80,0)}
                {\chi_{\alpha,\beta,\gamma}}{M_\gamma}{}
            \TNStackedMPOProduct{ftpRightB}{(2.00,0)}
                {\chi_{\alpha,\beta,\gamma}}{M_\gamma}{}
            \TNStackedMPOProduct{ftpRightC}{(3.20,0)}
                {\chi_{\alpha,\beta,\gamma}}{M_\gamma}{}
            \TNConnectVirtual{ftpRightAUpperEast}{ftpRightBUpperWest}
            \TNConnectVirtual{ftpRightBUpperEast}{ftpRightCUpperWest}
            \TNConnectVirtual{ftpRightALowerEast}{ftpRightBLowerWest}
            \TNConnectVirtual{ftpRightBLowerEast}{ftpRightCLowerWest}
            \TNTraceVirtualAbove{ftpRightAUpperWest}{ftpRightCUpperEast}
            \TNTraceVirtualBelow{ftpRightALowerWest}{ftpRightCLowerEast}
            \node at (2.00,-1.45) {$L\text{ sites}$};
    \end{TNDiagram}%
}

% Recursive contraction for the structure operator Q in CPSV16, lines
% 999--1008.  Each X has the four ordered indices k, alpha, beta, gamma;
% the lower wires retain earlier labels while the recursion proceeds.
\newcommand{\TNMPDORecursiveStructureOperator}{%
    \begin{TNDiagram}[normal]
        \path[use as bounding box] (-5.75,-1.55) rectangle (5.15,1.55);
        \node at (-5.25,0) {$Q=$};
        \TNMPOSite{rqOne}{(-4.15,0)}{X}
        \TNMPOSite{rqTwo}{(-2.70,0)}{X}
        \TNMPOSite{rqThree}{(-1.25,0)}{X}
        \TNConnectVirtual{rqOneEast}{rqTwoWest}
        \TNConnectVirtual{rqTwoEast}{rqThreeWest}
        \TNOpenVirtualWest{rqOneWest}
        \TNOpenPhysicalNorth{rqOneKet}
        \TNOpenPhysicalSouth{rqOneBra}
        \TNOpenPhysicalNorth{rqTwoKet}
        \TNOpenPhysicalSouth{rqTwoBra}
        \TNOpenPhysicalNorth{rqThreeKet}
        \TNOpenPhysicalSouth{rqThreeBra}
        \node at (0,0) {$\cdots$};
        \TNMPOSite{rqLast}{(1.25,0)}{X}
        \TNCompactTraceCell{rqTrace}{(3.35,0)}{M_\gamma}{1}{}
        \TNConnectVirtual{rqLastEast}{rqTraceInsertionW}
        \TNOpenPhysicalNorth{rqLastKet}
        \TNOpenPhysicalSouth{rqLastBra}
        \node[above] at (rqTwoKetOpen) {$k$};
        \node[below] at (rqTwoBraOpen) {$\alpha,\beta,\gamma$};
        \node[below] at (rqTraceCentre |- 0,-1.25) {$\tr(M_\gamma)$};
    \end{TNDiagram}%
}

% Fusion isometry of the renormalization fixed-point characterization: two
% contracted copies of the tensor equal one copy with the isometry merging
% the summed index into the physical pair.

% Appendix B basic-vector geometry.  This is the vector counterpart of the
% source figure II_RFP_MPS.png: each U acts on the two virtual halves at one
% physical site, while the two highlighted rank-one projectors act on the
% separate bonds shared by adjacent sites.
\newcommand{\TNAppendixBAdjacentBondProjectors}{%
    \begin{TNDiagram}[normal]
        \TNFusionMap{rfpUzero}{(-2.40,0.85)}{U_0}
        \TNFusionMap{rfpUone}{(0,0.85)}{U_1}
        \TNFusionMap{rfpUtwo}{(2.40,0.85)}{U_2}
        \TNMPSSite{rfpPhiZero}{(-1.20,-0.45)}{\varphi_{01}}
        \TNMPSSite{rfpPhiOne}{(1.20,-0.45)}{\varphi_{12}}
        \TNConnectVirtualVH{rfpUzeroFactorTwo}{rfpPhiZeroWest}
        \TNConnectVirtualVH{rfpUoneFactorOne}{rfpPhiZeroEast}
        \TNConnectVirtualVH{rfpUoneFactorTwo}{rfpPhiOneWest}
        \TNConnectVirtualVH{rfpUtwoFactorOne}{rfpPhiOneEast}
        \TNOpenVirtualNorth{rfpUzeroFactorOne}
        \TNOpenVirtualNorth{rfpUtwoFactorTwo}
        \TNOpenVirtualSouth{rfpUzeroCombined}
        \TNOpenVirtualSouth{rfpUoneCombined}
        \TNOpenVirtualSouth{rfpUtwoCombined}
        \node[tn grouping region, fit=(rfpPhiZero)] {};
        \node[tn grouping region, fit=(rfpPhiOne)] {};
        \node[red, below] at (rfpPhiZero) {$P_{01}$};
        \node[red, below] at (rfpPhiOne) {$P_{12}$};
    \end{TNDiagram}%
}

% Isometric transport of the two virtual bond projectors to the adjacent
% physical two-site basic-support projectors.  The upper row retains the
% source's three U tensors and diagonal bond vectors; the upper red brackets
% record the two overlapping physical supports.
\newcommand{\TNAppendixBPhysicalSupportTransport}{%
    \begin{TNDiagram}[normal]
        \TNFusionMap{rfpTzero}{(-2.40,0.65)}{U_0}
        \TNFusionMap{rfpTone}{(0,0.65)}{U_1}
        \TNFusionMap{rfpTtwo}{(2.40,0.65)}{U_2}
        \TNMPSSite{rfpTPhiZero}{(-1.20,-0.65)}{\varphi_{01}}
        \TNMPSSite{rfpTPhiOne}{(1.20,-0.65)}{\varphi_{12}}
        \TNConnectVirtualVH{rfpTzeroFactorTwo}{rfpTPhiZeroWest}
        \TNConnectVirtualVH{rfpToneFactorOne}{rfpTPhiZeroEast}
        \TNConnectVirtualVH{rfpToneFactorTwo}{rfpTPhiOneWest}
        \TNConnectVirtualVH{rfpTtwoFactorOne}{rfpTPhiOneEast}
        \TNOpenVirtualSouth{rfpTzeroCombined}
        \TNOpenVirtualSouth{rfpToneCombined}
        \TNOpenVirtualSouth{rfpTtwoCombined}
        \node[tn grouping region, fit=(rfpTzero)(rfpTone),
            label=above:{$(P_2)_{01}$}] {};
        \node[tn grouping region, fit=(rfpTone)(rfpTtwo),
            label=above:{$(P_2)_{12}$}] {};
        \node[red, below] at (rfpTPhiZero) {$P_{01}$};
        \node[red, below] at (rfpTPhiOne) {$P_{12}$};
    \end{TNDiagram}%
}

% Cyclic transport of the local three-site Appendix B commutator.  The left
% panel shows one three-site window inside an N-site ring; the right panel is
% the source AXB window with its two overlapping length-two interactions.
\newcommand{\TNAppendixBChainTransport}{%
    \begin{TNDiagram}[normal]
        \TNMPSSite{ctA}{(-2.20,0)}{A}
        \TNMPSSite{ctX}{(-1.10,1.15)}{X}
        \TNMPSSite{ctB}{(0.25,1.15)}{B}
        \TNMPSSite{ctC}{(1.35,0)}{}
        \TNMPSSite{ctD}{(0.25,-1.15)}{}
        \TNMPSSite{ctE}{(-1.10,-1.15)}{}
        \TNConnectVirtual{ctAEast}{ctXWest}
        \TNConnectVirtual{ctXEast}{ctBWest}
        \TNConnectVirtual{ctBEast}{ctCWest}
        \TNConnectVirtual{ctCEast}{ctDWest}
        \TNConnectVirtual{ctDEast}{ctEWest}
        \TNConnectVirtual{ctEEast}{ctAWest}
        \node[red, above left] at (ctX) {$h_i$};
        \node[blue, above right] at (ctX) {$h_{i+1}$};
        \node at (-0.45,-1.75) {$N>2$};
        \TNMorphismArrow{ctTransport}{(2.00,0)}{(3.15,0)}{R_{i,\tau}^{(3)}}
        \TNMPSSite{ctLA}{(3.80,0)}{A}
        \TNMPSSite{ctLX}{(5.20,0)}{X}
        \TNMPSSite{ctLB}{(6.60,0)}{B}
        \TNConnectVirtual{ctLAEast}{ctLXWest}
        \TNConnectVirtual{ctLXEast}{ctLBWest}
        \TNOpenVirtualWest{ctLAWest}
        \TNOpenVirtualEast{ctLBEast}
        \node[red, below] at ($(ctLA)!0.5!(ctLX)$) {$h_0^{(3)}$};
        \node[blue, below] at ($(ctLX)!0.5!(ctLB)$) {$h_1^{(3)}$};
    \end{TNDiagram}%
}

\newcommand{\TNRFPKrausIsometry}{%
    \begin{tikzpicture}[tn picture]
        \TN@mpssite{rfkA}{(0,0)}{}
        \TN@mpssite{rfkB}{(1.20,0)}{}
        \TN@vopenport{rfkAW}{-\TN@virtualleglen,0}
        \TNConnectVirtual{rfkAE}{rfkBW}
        \TN@vopenport{rfkBE}{\TN@virtualleglen,0}
        \TN@popenport{rfkAN}{0,\TN@physicalleglen}
        \TN@popenport{rfkBN}{0,\TN@physicalleglen}
        \TN@label{rfkANOpen}{(0,0.24)}{i_1}
        \TN@label{rfkBNOpen}{(0,0.24)}{i_2}
        \node[anchor=east] at (-0.14,-0.28) {$A$};
        \node[anchor=west] at (1.34,-0.28) {$A$};
        \node at (2.50,0) {$=$};
        \node at (3.25,0) {$\displaystyle\sum_j$};
        \TN@mpssite{rfkC}{(4.30,0)}{}
        \TN@openhorizontalports{rfkC}
        \node[anchor=west] at (4.44,-0.28) {$A$};
        \TN@physicalsplitmap[inner xsep=10pt, inner ysep=2pt]
            {rfkV}{(4.30,0.90)}{V}
        \TNConnectPhysical{rfkCN}{rfkVTrunk}
        \node[anchor=west] at (4.40,0.46) {$j$};
        \TN@popenport{rfkVLeft}{0,0.34}
        \TN@popenport{rfkVRight}{0,0.34}
        \node at (4.10,1.62) {$i_1$};
        \node at (4.50,1.62) {$i_2$};
    \end{tikzpicture}%
}

% Reverse fusion identity: the adjoint isometry applied to the two physical
% legs of the contracted pair recovers the single tensor.
\newcommand{\TNRFPKrausIsometryReverse}{%
    \begin{tikzpicture}[tn picture]
        \TN@mpssite{rfrA}{(0,0)}{}
        \TN@mpssite{rfrB}{(1.20,0)}{}
        \TN@vopenport{rfrAW}{-\TN@virtualleglen,0}
        \TNConnectVirtual{rfrAE}{rfrBW}
        \TN@vopenport{rfrBE}{\TN@virtualleglen,0}
        \node[anchor=east] at (-0.14,-0.28) {$A$};
        \node[anchor=west] at (1.34,-0.28) {$A$};
        \TN@physicalmergemap[minimum width=1.62cm]
            {rfrV}{(0.60,0.90)}{V^\dagger}
        \TNConnectPhysical{rfrAN}{rfrVLeft}
        \TNConnectPhysical{rfrBN}{rfrVRight}
        \TN@popenport{rfrVTrunk}{0,0.34}
        \node[anchor=west] at (0.70,1.36) {$j$};
        \node at (2.50,0) {$=$};
        \TN@mpssite{rfrC}{(3.55,0)}{}
        \TN@openhorizontalports{rfrC}
        \TN@popenport{rfrCN}{0,\TN@physicalleglen}
        \TN@label{rfrCNOpen}{(0,0.24)}{j}
        \node[anchor=west] at (3.69,-0.28) {$A$};
    \end{tikzpicture}%
}

% Isometry canonical form of a normal renormalization fixed point: the
% virtual chain X, the diagonal weight, the isometry emitting the physical
% leg, and X^{-1}.
\newcommand{\TNRFPIsometryCanonicalForm}{%
    \begin{tikzpicture}[tn picture]
        \TNTensor{icfA}{(0,0)}
        \TN@vleg{icfA}{west}{-0.55,0}
        \TN@vleg{icfA}{east}{0.55,0}
        \TN@pleg{icfA}{north}{0,0.46}
        \TN@label{icfA.north}{(0,0.70)}{i}
        \node[anchor=east] at (-0.14,-0.28) {$A$};
        \node at (1.25,0) {$=$};
        \TN@insertion[label=above:$X$]{icfX}{(2.20,0)}
        \TN@insertion[label=above:$\sqrt\Lambda$]{icfL}{(3.20,0)}
        \TN@bondpairmap{icfU}{(4.30,0.75)}{U}
        \TN@insertion[label=above:$X^{-1}$]{icfXi}{(5.40,0)}
        \TN@vport{icfXW}{icfX}{west}
        \TN@vport{icfXE}{icfX}{east}
        \TN@vport{icfLW}{icfL}{west}
        \TN@vport{icfLE}{icfL}{east}
        \TN@vport{icfXiW}{icfXi}{west}
        \TN@vport{icfXiE}{icfXi}{east}
        \TN@vopenport{icfXW}{-0.55,0}
        \TNConnectVirtual{icfXE}{icfLW}
        \TNConnectVirtualHV{icfLE}{icfULeft}
        \TNConnectVirtualVH{icfURight}{icfXiW}
        \TN@vopenport{icfXiE}{0.55,0}
        \TN@popenport{icfUPhysical}{0,0.34}
        \node at ($(icfUPhysicalOpen)+(0,0.22)$) {$i$};
    \end{tikzpicture}%
}

% Per-block isometry canonical form of a renormalization fixed point in
% canonical form: the same chain threaded by the block-index lines j and q,
% with the coefficient tensor mu on the two index lines.
\newcommand{\TNRFPIsometryCanonicalFormBlocks}{%
    \begin{tikzpicture}[tn picture]
        \TNTensor{icbA}{(0,0)}
        \TN@vleg{icbA}{west}{-0.55,0}
        \TN@vleg{icbA}{east}{0.55,0}
        \TN@pleg{icbA}{north}{0,0.46}
        \TN@label{icbA.north}{(0,0.70)}{i}
        \node[anchor=east] at (-0.14,-0.28) {$A$};
        \node at (1.25,0) {$=$};
        \TN@sectorgauge{icbX}{(2.20,0)}{X}
        \TN@insertion[label=above:$\sqrt\Lambda$]{icbL}{(3.20,0)}
        \TN@bondpairmap{icbU}{(4.30,0.75)}{U}
        \TN@inverseSectorgauge{icbXi}{(5.40,0)}{X^{-1}}
        \TN@vport{icbLW}{icbL}{west}
        \TN@vport{icbLE}{icbL}{east}
        \TN@vopenport{icbXW}{-0.55,0}
        \TNConnectVirtual{icbXE}{icbLW}
        \TNConnectVirtualHV{icbLE}{icbULeft}
        \TNConnectVirtualVH{icbURight}{icbXiW}
        \TN@vopenport{icbXiE}{0.55,0}
        \TN@popenport{icbUPhysical}{0,0.34}
        \node at ($(icbUPhysicalOpen)+(0,0.22)$) {$i$};
        \TN@vbus{icbJ}{(1.65,-0.62)}{(5.95,-0.62)}{j}
        \TN@vparallelbus{icbQ}{icbJ}{-\TN@buspitch}{q}
        \TN@vport{icbLBus}{icbL}{south}
        \TN@vbusprojecttap{icbJ}{icbJX}{icbXBlock}
        \TN@vbusprojecttap{icbQ}{icbQX}{icbXCopy}
        \TNConnectVirtual{icbXBlock}{icbJX}
        \TNConnectVirtual{icbXCopy}{icbQX}
        \TN@vbusprojecttap{icbJ}{icbJL}{icbLBus}
        \TNConnectVirtual{icbLBus}{icbJL}
        \TN@vterminal{icbUBlock}{(icbULeft |- icbL)}
        \TN@anonymousjunction{(icbUBlock)}
        \TN@vbusprojecttap{icbJ}{icbJU}{icbUBlock}
        \TNConnectVirtual{icbUBlock}{icbJU}
        \TN@vbusprojecttap{icbJ}{icbJXi}{icbXiBlock}
        \TN@vbusprojecttap{icbQ}{icbQXi}{icbXiCopy}
        \TNConnectVirtual{icbXiBlock}{icbJXi}
        \TNConnectVirtual{icbXiCopy}{icbQXi}
        \TNComponent{icbMu}{(2.70,-1.60)}{\mu}
        \TN@vnorthport{icbMuNLeft}{icbMu}{0.15}
        \TN@vnorthport{icbMuNRight}{icbMu}{0.85}
        \TN@vbusprojecttap{icbJ}{icbJMu}{icbMuNLeft}
        \TN@vbusprojecttap{icbQ}{icbQMu}{icbMuNRight}
        \TNConnectVirtual{icbMuNLeft}{icbJMu}
        \TNConnectVirtual{icbMuNRight}{icbQMu}
    \end{tikzpicture}%
}

% Two-site blocking of a simple MPO tensor: two copies contracted along one
% virtual bond, with the two ket legs and the two bra legs grouped into the
% blocked physical indices (dashed box).
\newcommand{\TNMPDOTwoSiteBlocking}{%
    \begin{TNEquationRow}[compact]
        \TNTerm{%
            \TNMPOSite{tsbM}{(0,0)}{M}
            \TNOpenVirtualWest{tsbMWest}
            \TNOpenVirtualEast{tsbMEast}
            \TNOpenPhysicalNorth{tsbMKet}
            \TNOpenPhysicalSouth{tsbMBra}}
        \TNRelation{=}
        \TNTerm{%
            \TNMPOSite{tsbKa}{(0,0)}{K}
            \TNMPOSite{tsbKb}{(1.25,0)}{K}
            \TNOpenVirtualWest{tsbKaWest}
            \TNConnectVirtual{tsbKaEast}{tsbKbWest}
            \TNOpenVirtualEast{tsbKbEast}
            \TNOpenPhysicalNorth{tsbKaKet}
            \TNOpenPhysicalSouth{tsbKaBra}
            \TNOpenPhysicalNorth{tsbKbKet}
            \TNOpenPhysicalSouth{tsbKbBra}}
    \end{TNEquationRow}%
}

% Vertical product of two distinct basis tensors, contracted through the
% shared physical index.
\newcommand{\TNMPDOBNTVerticalProduct}{%
    \begin{TNEquationRow}[compact]
        \TNTerm{%
            \TNStackedMPOProduct{vp}{(0,0)}{\mathcal K_y}{\mathcal K_x}{}
            \TNOpenVirtualWest{vpUpperWest}
            \TNOpenVirtualEast{vpUpperEast}
            \TNOpenVirtualWest{vpLowerWest}
            \TNOpenVirtualEast{vpLowerEast}
            \TNOpenPhysicalNorth{vpUpperKet}
            \TNOpenPhysicalSouth{vpLowerBra}}
        \TNRelation{=}
        \ensuremath{0\qquad(x\ne y)}
    \end{TNEquationRow}%
}

% The two block closure maps: a boundary insertion on the virtual bond of a
% block tensor, with the virtual loop closed by the trace, gives an operator
% on the open physical legs.
\newcommand{\TNMPDOBlockClosureMap}{%
    \begin{TNEquationRow}[compact]
        \ensuremath{M_\alpha(X)=}
        \TNTerm{\TNCompactTraceCell{bcm}{(0,0)}{M_\alpha}{X}{}}
        \qquad
        \ensuremath{A_\alpha(X)=}
        \TNTerm{\TNCompactTraceCell{bca}{(0,0)}{A_\alpha}{X}{}}
    \end{TNEquationRow}%
}

% The BNT-label operator O_L(M_alpha): L vertically read copies of M_alpha
% are multiplied along the vertical index and closed by the trace, while the
% horizontal legs remain open.  This follows source figure MPDO_O_L_A_.png.
\newcommand{\TNMPDOBNTOperatorTraceClosure}{%
    \begin{TNEquationRow}[compact]
        \ensuremath{O_L(M_\alpha)=}
        \TNTerm{\TNVerticalWord{botc}{(0,0)}{M_\alpha}{closed}}
    \end{TNEquationRow}%
}

% First-site insertion agreement on the closed word: inserting Y on the
% physical leg of the first site of the matrix product vector gives the
% same vector as inserting Z.
\newcommand{\TNMPDOFirstSiteInsertionHypothesis}{%
    \begin{tikzpicture}[tn picture]
        \begin{scope}
            \TN@declareWordTensor{fihL}{(0,0)}
            \TN@declareWordTensor{fihM}{(1.05,0)}
            \TN@declareWordTensor{fihR}{(2.60,0)}
            \TNConnectVirtual{fihLE}{fihMW}
            \TN@vopenport{fihME}{0.45,0}
            \TN@vopenport{fihRW}{-0.45,0}
            \TN@vtracepath{(fihLW) -- (-0.42,0) -- (-0.42,-0.44)
                -- (3.02,-0.44) -- (3.02,0) -- (fihRE)}
            \TN@popenport{fihMN}{0,0.42}
            \TN@popenport{fihRN}{0,0.42}
            \TN@insertion[label=left:$Y$]{fihY}{(0,0.46)}
            \TN@pport{fihYS}{fihY}{south}
            \TN@pport{fihYN}{fihY}{north}
            \TNConnectPhysical{fihLN}{fihYS}
            \TN@popenport{fihYN}{0,0.40}
            \node at (1.85,0) {$\cdots$};
        \end{scope}
        \node at (3.60,0) {$=$};
        \begin{scope}[xshift=4.20cm]
            \TN@declareWordTensor{fizL}{(0,0)}
            \TN@declareWordTensor{fizM}{(1.05,0)}
            \TN@declareWordTensor{fizR}{(2.60,0)}
            \TNConnectVirtual{fizLE}{fizMW}
            \TN@vopenport{fizME}{0.45,0}
            \TN@vopenport{fizRW}{-0.45,0}
            \TN@vtracepath{(fizLW) -- (-0.42,0) -- (-0.42,-0.44)
                -- (3.02,-0.44) -- (3.02,0) -- (fizRE)}
            \TN@popenport{fizMN}{0,0.42}
            \TN@popenport{fizRN}{0,0.42}
            \TN@insertion[label=left:$Z$]{fizZ}{(0,0.46)}
            \TN@pport{fizZS}{fizZ}{south}
            \TN@pport{fizZN}{fizZ}{north}
            \TNConnectPhysical{fizLN}{fizZS}
            \TN@popenport{fizZN}{0,0.40}
            \node at (1.85,0) {$\cdots$};
        \end{scope}
    \end{tikzpicture}%
}

% Blockwise conclusion of the first-site agreement: the one-site insertions
% of Y and Z agree on each basis tensor of the canonical form.
\newcommand{\TNMPDOFirstSiteInsertionBlockwise}{%
    \begin{tikzpicture}[tn picture]
        \TNTensor{fibA}{(0,0)}
        \TN@vleg{fibA}{west}{-0.60,0}
        \TN@vleg{fibA}{east}{0.60,0}
        \TN@pleg{fibA}{north}{0,0.86}
        \TN@insertion[label=left:$Y$]{fibY}{(0,0.46)}
        \node at (0,-0.32) {$A_k$};
        \node at (1.45,0) {$=$};
        \TNTensor{fibB}{(2.90,0)}
        \TN@vleg{fibB}{west}{-0.60,0}
        \TN@vleg{fibB}{east}{0.60,0}
        \TN@pleg{fibB}{north}{0,0.86}
        \TN@insertion[label=left:$Z$]{fibZ}{(2.90,0.46)}
        \node at (2.90,-0.32) {$A_k$};
    \end{tikzpicture}%
}

% Trace pairing of a one-site insertion against a boundary operator: the
% virtual loop closes the inserted site against W, and the pairing takes the
% same value for the Y-insertion and the Z-insertion.
\newcommand{\TNMPDOInsertionTracePairing}{%
    \begin{tikzpicture}[tn picture]
        \begin{scope}
            \TNComponent{itpW}{(0,0)}{W}
            \TNTensor{itpA}{(1.15,0)}
            \TN@vconnect{itpW}{east}{itpA}{west}
            \TN@pleg{itpA}{north}{0,0.86}
            \TN@insertion[label=left:$Y$]{itpY}{(1.15,0.46)}
            \TN@vtracepath{(itpW.west) -- ++(-0.35,0) -- (-0.62,-0.80)
            -- (1.72,-0.80) -- (1.72,0) -- (itpA.east)}
            \node at (0.55,-1.03) {$\tr$};
        \end{scope}
        \node at (2.70,0) {$=$};
        \begin{scope}[xshift=4.05cm]
            \TNComponent{itpWz}{(0,0)}{W}
            \TNTensor{itpAz}{(1.15,0)}
            \TN@vconnect{itpWz}{east}{itpAz}{west}
            \TN@pleg{itpAz}{north}{0,0.86}
            \TN@insertion[label=left:$Z$]{itpZ}{(1.15,0.46)}
            \TN@vtracepath{(itpWz.west) -- ++(-0.35,0) -- (-0.62,-0.80)
            -- (1.72,-0.80) -- (1.72,0) -- (itpAz.east)}
            \node at (0.55,-1.03) {$\tr$};
        \end{scope}
    \end{tikzpicture}%
}

% Horizontal canonical form of the doubled tensor: the gauge X, one basis
% tensor carrying the ket and bra legs, and X^{-1}, threaded by the
% block-index lines j and q; the weights are absorbed into the basis
% tensors, so the basis tensor taps only the block line j while the gauge
% factors tap both lines.
\newcommand{\TNMPDOHorizontalCanonicalForm}{%
    \begin{tikzpicture}[tn picture]
        \TNTensor{hcfK}{(0,0)}
        \TN@vleg{hcfK}{west}{-0.58,0}
        \TN@vleg{hcfK}{east}{0.58,0}
        \TN@pleg{hcfK}{north}{0,0.50}
        \TN@pleg{hcfK}{south}{0,-0.50}
        \node[anchor=east] at (-0.14,-0.28) {$K$};
        \node at (1.30,0) {$=$};
        \TN@sectorgauge{hcfX}{(2.55,0)}{X}
        \TNTensor{hcfC}{(3.75,0)}
        \TN@inverseSectorgauge{hcfXi}{(4.95,0)}{X^{-1}}
        \TN@vleg{hcfX}{west}{-0.55,0}
        \TN@vconnect{hcfX}{east}{hcfC}{west}
        \TN@vconnect{hcfC}{east}{hcfXi}{west}
        \TN@vleg{hcfXi}{east}{0.55,0}
        \TN@pleg{hcfC}{north}{0,0.50}
        \TN@pleg{hcfC}{south}{0,-0.50}
        \node[anchor=west] at (3.89,0.32) {$\mathcal K$};
        \TN@vbus{hcfJ}
            {($(hcfX.center)+(-0.60,-\TN@busoffset)$)}
            {($(hcfXi.center)+(0.60,-\TN@busoffset)$)}{j}
        \TN@vparallelbus{hcfQ}{hcfJ}{-\TN@buspitch}{q}
        \TN@vbusprojecttap{hcfJ}{hcfJX}{hcfXBlock}
        \TN@vbusprojecttap{hcfQ}{hcfQX}{hcfXCopy}
        \TNConnectVirtual{hcfXBlock}{hcfJX}
        \TNConnectVirtual{hcfXCopy}{hcfQX}
        \TN@vport{hcfCBlock}{hcfC}{-45}
        \TN@vbusprojecttap{hcfJ}{hcfJC}{hcfCBlock}
        \TNConnectVirtual{hcfCBlock}{hcfJC}
        \TN@vbusprojecttap{hcfJ}{hcfJXi}{hcfXiBlock}
        \TN@vbusprojecttap{hcfQ}{hcfQXi}{hcfXiCopy}
        \TNConnectVirtual{hcfXiBlock}{hcfJXi}
        \TNConnectVirtual{hcfXiCopy}{hcfQXi}
    \end{tikzpicture}%
}

% Inverse tensor of a block-injective canonical form: contracting the
% inverse tensor with the doubled physical index of the canonical form
% passes the virtual legs of the gauge factors to the virtual legs of the
% inverse tensor and ties the block index of the inverse tensor to the
% block line.
\newcommand{\TNMPDOBlockInjectiveInverse}{%
    \begin{tikzpicture}[tn picture]
        \TNTensor{bijI}{(0,0.55)}
        \TNTensor{bijK}{(0,-0.55)}
        \TN@vleg{bijI}{west}{-0.55,0}
        \TN@vleg{bijI}{east}{0.55,0}
        \TN@vleg{bijK}{west}{-0.55,0}
        \TN@vleg{bijK}{east}{0.55,0}
        \TN@pconnect{bijI}{south}{bijK}{north}
        \TN@vleg{bijI}{north}{0,0.42}
        \node at (-0.58,0.90) {$K^{-1}$};
        \node at (-0.42,-0.90) {$K$};
        \node at (1.15,0) {$=$};
        \begin{scope}[xshift=2.30cm]
            \TN@sectorgauge{bijX}{(0.55,0)}{X}
            \TNTensor{bijC}{(1.75,0)}
            \TN@inverseSectorgauge{bijXi}{(2.95,0)}{X^{-1}}
            \TN@vleg{bijX}{west}{-0.50,0}
            \TN@vconnect{bijX}{east}{bijC}{west}
            \TN@vconnect{bijC}{east}{bijXi}{west}
            \TN@vleg{bijXi}{east}{0.50,0}
            \node[anchor=west] at (1.89,-0.30) {$\mathcal K$};
            \TNTensor{bijIt}{(1.75,1.00)}
            \TN@vleg{bijIt}{west}{-0.50,0}
            \TN@vleg{bijIt}{east}{0.50,0}
            \TN@pconnect{bijIt}{south}{bijC}{north}
            \TN@vleg{bijIt}{north}{0,0.42}
            \node at (1.02,1.32) {$K^{-1}$};
            \TN@vbus{bijJ}
                {($(bijX.center)+(-0.55,-\TN@busoffset)$)}
                {($(bijXi.center)+(0.55,-\TN@busoffset)$)}{j}
            \TN@vparallelbus{bijQ}{bijJ}{-\TN@buspitch}{q}
            \TN@vport{bijCBlock}{bijC}{south}
            \TN@vbusprojecttap{bijJ}{bijJX}{bijXBlock}
            \TN@vbusprojecttap{bijQ}{bijQX}{bijXCopy}
            \TNConnectVirtual{bijXBlock}{bijJX}
            \TNConnectVirtual{bijXCopy}{bijQX}
            \TN@vbusprojecttap{bijJ}{bijJC}{bijCBlock}
            \TNConnectVirtual{bijCBlock}{bijJC}
            \TN@vbusprojecttap{bijJ}{bijJXi}{bijXiBlock}
            \TN@vbusprojecttap{bijQ}{bijQXi}{bijXiCopy}
            \TNConnectVirtual{bijXiBlock}{bijJXi}
            \TNConnectVirtual{bijXiCopy}{bijQXi}
        \end{scope}
        \node at (6.25,0) {$=$};
        \begin{scope}[xshift=6.70cm]
            \TN@sectorgauge{bijPX}{(0.55,0)}{X}
            \TN@inverseSectorgauge{bijPXi}{(2.95,0)}{X^{-1}}
            \TN@vleg{bijPX}{west}{-0.50,0}
            \TN@vleg{bijPXi}{east}{0.50,0}
            \TN@vpath[rounded corners=3pt]{(bijPX.east) -- (1.05,0) -- (1.05,1.00) -- (1.35,1.00)}
            \TN@vpath[rounded corners=3pt]{(bijPXi.west) -- (2.45,0) -- (2.45,1.00) -- (2.15,1.00)}
            \TN@vbus{bijPJ}
                {($(bijPX.center)+(-0.55,-\TN@busoffset)$)}
                {($(bijPXi.center)+(0.55,-\TN@busoffset)$)}{}
            \TN@vparallelbus{bijPQ}{bijPJ}{-\TN@buspitch}{}
            \TN@vbusprojecttap{bijPJ}{bijPJX}{bijPXBlock}
            \TN@vbusprojecttap{bijPQ}{bijPQX}{bijPXCopy}
            \TNConnectVirtual{bijPXBlock}{bijPJX}
            \TNConnectVirtual{bijPXCopy}{bijPQX}
            \TN@vbusprojecttap{bijPJ}{bijPJXi}{bijPXiBlock}
            \TN@vbusprojecttap{bijPQ}{bijPQXi}{bijPXiCopy}
            \TNConnectVirtual{bijPXiBlock}{bijPJXi}
            \TNConnectVirtual{bijPXiCopy}{bijPQXi}
            \TN@vterminal{bijPBlockTop}{(1.75,1.50)}
            \TN@vbusprojecttap{bijPJ}{bijPJBlock}{bijPBlockTop}
            \TNConnectVirtual{bijPBlockTop}{bijPJBlock}
        \end{scope}
    \end{tikzpicture}%
}

% Recovery of the tensor from the inverse map: contracting the inverse
% tensor with the doubled physical index and closing its virtual legs
% against a further copy of the tensor returns the tensor itself; the
% doubled physical index is drawn as a single thick leg.
\newcommand{\TNMPDOInverseRecovery}{%
    \begin{tikzpicture}[tn picture]
        \TNTensor{irT}{(0,1.05)}
        \TNTensor{irI}{(0,0)}
        \TNTensor{irB}{(0,-1.05)}
        \TN@pleg{irT}{north}{0,0.45}
        \TN@vpath[rounded corners=3pt]{(irI.west) -- ++(-0.74,0) -- (-0.82,1.05) -- (irT.west)}
        \TN@vpath[rounded corners=3pt]{(irI.east) -- ++(0.74,0) -- (0.82,1.05) -- (irT.east)}
        \TN@pconnect{irI}{south}{irB}{north}
        \TN@vleg{irB}{west}{-0.60,0}
        \TN@vleg{irB}{east}{0.60,0}
        \node[anchor=west] at (0.12,1.33) {$K$};
        \node[anchor=west] at (0.12,0.30) {$K^{-1}$};
        \node[anchor=east] at (-0.14,-1.33) {$K$};
        \node at (1.85,0) {$=$};
        \TNTensor{irR}{(3.05,0)}
        \TN@vleg{irR}{west}{-0.60,0}
        \TN@vleg{irR}{east}{0.60,0}
        \TN@pleg{irR}{north}{0,0.45}
        \node[anchor=east] at (2.91,-0.28) {$K$};
    \end{tikzpicture}%
}

% Projector absorption for a basis tensor: expanding the identity as the
% sum of the orthogonal projectors on the ket and bra legs leaves only the
% diagonal term, so the tensor absorbs its own projector on either side.
\newcommand{\TNMPDOProjectorAbsorption}{%
    \begin{TNEquationRow}[compact]
        \ensuremath{
            \mathcal K_i
            =\sum_{j,j'}P_j\,\mathcal K_i\,P_{j'}
            =P_i\,\mathcal K_i
            =\mathcal K_i\,P_i}
    \end{TNEquationRow}%
}

\newcommand{\TNGaugeConjugation}[3]{%
    \begin{tikzpicture}[tn picture]
        \TN@insertion[label=above:$#1$]{tnXl}{(-1.10,0)}
        \TN@tensor{tnA}{(0,0)}
        \TN@insertion[label=above:$#3$]{tnXr}{(1.10,0)}
        \TN@vleg{tnXl}{west}{-0.55,0}
        \TN@vconnect{tnXl}{east}{tnA}{west}
        \TN@vconnect{tnA}{east}{tnXr}{west}
        \TN@vleg{tnXr}{east}{0.55,0}
        \TN@pleg{tnA}{north}{0,\TN@physleglen}
        \TN@uplabel{tnA}{\TN@physlabeloff}{#2}
    \end{tikzpicture}%
}

\newcommand{\TNPhysicalRealization}[2]{%
    \begin{tikzpicture}[tn picture]
        \TN@tensor{tnA1}{(0,0)}
        \TN@vleg{tnA1}{west}{-0.65,0}
        \TN@pleg{tnA1}{north}{0,\TN@physleglen}
        \TN@insertion[label=above:$#1$]{tnX}{(0.78,0)}
        \TN@vconnect{tnA1}{east}{tnX}{west}
        \TN@vleg{tnX}{east}{0.22,0}
        \node at (1.78,0) {$=$};
        \TN@tensor{tnA2}{(2.65,0)}
        \TN@vleg{tnA2}{west}{-0.65,0}
        \TN@vleg{tnA2}{east}{0.65,0}
        \TN@insertion[label=left:$#2$]{tnO}{(2.65,0.74)}
        \TN@pconnect{tnA2}{north}{tnO}{south}
    \end{tikzpicture}%
}

% Fix (#401): use a dedicated physical-leg operator insertion for
% general linear twists, and keep permutation twists for index relabellings.
\newcommand{\TNLinearTwist}[2]{%
    \begin{tikzpicture}[tn picture]
        \TN@tensor{tnA}{(0,0)}
        \TN@vleg{tnA}{west}{-0.65,0}
        \TN@vleg{tnA}{east}{0.65,0}
        \TN@insertion[label=left:$#1$]{tnu}{(0,0.64)}
        \TN@pconnect{tnA}{north}{tnu}{south}
        \node at ($(tnu.north)+(0,0.18)$) {$#2$};
    \end{tikzpicture}%
}

% Fix (#401): make the relabelling arrow explicit so multi-step
% permutation diagrams can distinguish \sigma_h from \sigma_g.
\newcommand{\TNPermutationTwistLabeled}[3]{%
    \begin{tikzpicture}[tn picture]
        \TN@tensor{tnL}{(0,0)}
        \TN@vleg{tnL}{west}{-0.65,0}
        \TN@vleg{tnL}{east}{0.65,0}
        \TN@pleg{tnL}{north}{0,\TN@physleglen}
        \TN@uplabel{tnL}{\TN@physlabeloff}{#1}
        \node at (1.55,0) {$\xrightarrow{\ #3\ }$};
        \TN@tensor{tnR}{(3.30,0)}
        \TN@vleg{tnR}{west}{-0.65,0}
        \TN@vleg{tnR}{east}{0.65,0}
        \TN@pleg{tnR}{north}{0,\TN@physleglen}
        \TN@uplabel{tnR}{\TN@physlabeloff}{#2}
    \end{tikzpicture}%
}

\newcommand{\TNPermutationTwist}[2]{%
    \TNPermutationTwistLabeled{#1}{#2}{\sigma_g}%
}

\newcommand{\TNTwistedTransfer}[1]{%
    \begin{tikzpicture}[tn picture]
        \TN@tensor{tnUp}{(0,\TN@doublelayersep)}
        \TN@tensor{tnDn}{(0,-\TN@doublelayersep)}
        \TN@vleg{tnUp}{west}{-0.80,0}
        \TN@vleg{tnDn}{west}{-0.80,0}
        \TN@vleg{tnUp}{east}{0.80,0}
        \TN@vleg{tnDn}{east}{0.80,0}
        \TN@insertion[label=left:$#1$]{tnu}{(0,0)}
        \TN@pconnect{tnUp}{south}{tnu}{north}
        \TN@pconnect{tnu}{south}{tnDn}{north}
        \node at (-1.25,0) {$X$};
        \node at (1.52,0) {$\E_{#1}(X)$};
        \node at ($(tnUp.north)+(0,0.28)$) {$n$};
        \node at ($(tnDn.south)+(0,-0.28)$) {$n'$};
    \end{tikzpicture}%
}

\newcommand{\TNCondCOne}[3]{%
    \begin{tikzpicture}[tn picture]
        \TN@tensor{tnL}{(0,0)}
        \TN@vleg{tnL}{west}{-0.65,0}
        \TN@vleg{tnL}{east}{0.65,0}
        \TN@pleg{tnL}{north}{0,0.26}
        \TN@insertion[label=left:$#1$]{tnu}{(0,0.64)}
        \TN@pconnect{tnL}{north}{tnu}{south}
        \node at ($(tnu.north)+(0,0.18)$) {$#3$};
        \node at (1.65,0) {$=$};
        \TN@insertion[label=above:$#2$]{tnXl}{(3.15,0)}
        \TN@tensor{tnA}{(4.20,0)}
        \TN@insertion[label=above:$#2^\dagger$]{tnXr}{(5.25,0)}
        \TN@vleg{tnXl}{west}{-0.55,0}
        \TN@vconnect{tnXl}{east}{tnA}{west}
        \TN@vconnect{tnA}{east}{tnXr}{west}
        \TN@vleg{tnXr}{east}{0.55,0}
        \TN@pleg{tnA}{north}{0,\TN@physleglen}
        \TN@uplabel{tnA}{\TN@physlabeloff}{#3}
    \end{tikzpicture}%
}

\newcommand{\TNCondCTwo}[1]{%
    \begin{tikzpicture}[tn picture]
        \TN@insertion[label=above:$#1$]{tnVL}{(-1.15,0.48)}
        \TN@insertion[label=below:$#1^\dagger$]{tnVdL}{(-1.15,-0.48)}
        \TN@tensor{tnLUp}{(0,0.48)}
        \TN@tensor{tnLDn}{(0,-0.48)}
        \TN@vleg{tnVL}{west}{-0.45,0}
        \TN@vleg{tnVdL}{west}{-0.45,0}
        \TN@vconnect{tnVL}{east}{tnLUp}{west}
        \TN@vconnect{tnVdL}{east}{tnLDn}{west}
        \TN@vleg{tnLUp}{east}{0.55,0}
        \TN@vleg{tnLDn}{east}{0.55,0}
        \TN@pconnect{tnLUp}{south}{tnLDn}{north}
        \node at (1.55,0) {$=$};
        \TN@tensor{tnRUp}{(3.10,0.48)}
        \TN@tensor{tnRDn}{(3.10,-0.48)}
        \TN@insertion[label=above:$#1$]{tnVR}{(4.25,0.48)}
        \TN@insertion[label=below:$#1^\dagger$]{tnVdR}{(4.25,-0.48)}
        \TN@vleg{tnRUp}{west}{-0.55,0}
        \TN@vleg{tnRDn}{west}{-0.55,0}
        \TN@vconnect{tnRUp}{east}{tnVR}{west}
        \TN@vconnect{tnRDn}{east}{tnVdR}{west}
        \TN@vleg{tnVR}{east}{0.45,0}
        \TN@vleg{tnVdR}{east}{0.45,0}
        \TN@pconnect{tnRUp}{south}{tnRDn}{north}
    \end{tikzpicture}%
}

\newcommand{\TNStringOrderParameter}[2]{%
    \begin{tikzpicture}[tn picture]
        \TN@tensor{tnLam}{(-1.15,0)}
        \node at ($(tnLam.south)+(0,-0.28)$) {$\Lambda$};
        \TN@tensor{tnU1}{(0,\TN@doublelayersep)}
        \TN@tensor{tnD1}{(0,-\TN@doublelayersep)}
        \TN@tensor{tnU2}{(1.55,\TN@doublelayersep)}
        \TN@tensor{tnD2}{(1.55,-\TN@doublelayersep)}
        \TN@tensor{tnU3}{(4.30,\TN@doublelayersep)}
        \TN@tensor{tnD3}{(4.30,-\TN@doublelayersep)}
        \TN@insertion[label=left:$#1$]{tnuA}{(0,0)}
        \TN@insertion[label=left:$#1$]{tnuB}{(1.55,0)}
        \TN@insertion[label=left:$#1$]{tnuC}{(4.30,0)}
        \TN@vconnect{tnLam}{east}{tnU1}{west}
        \TN@vconnect{tnLam}{east}{tnD1}{west}
        \TN@vconnect{tnU1}{east}{tnU2}{west}
        \TN@vconnect{tnD1}{east}{tnD2}{west}
        \TN@vleg{tnU2}{east}{0.62,0}
        \TN@vleg{tnD2}{east}{0.62,0}
        \TN@vleg{tnU3}{west}{-0.62,0}
        \TN@vleg{tnD3}{west}{-0.62,0}
        \TN@vtraceright{tnU3}{tnD3}{0.70}
        \TN@pconnect{tnU1}{south}{tnuA}{north}
        \TN@pconnect{tnuA}{south}{tnD1}{north}
        \TN@pconnect{tnU2}{south}{tnuB}{north}
        \TN@pconnect{tnuB}{south}{tnD2}{north}
        \TN@pconnect{tnU3}{south}{tnuC}{north}
        \TN@pconnect{tnuC}{south}{tnD3}{north}
        \node at (2.92,\TN@doublelayersep) {$\cdots$};
        \node at (2.92,-\TN@doublelayersep) {$\cdots$};
        \node at (5.55,0) {$\operatorname{tr}$};
        \node at (2.92,-1.02) {$#2$};
    \end{tikzpicture}%
}

\newcommand{\TNInternalTraceInsertion}[3]{%
    \begin{tikzpicture}[tn picture]
        \TN@tensor{tnL}{(0,0)}
        \TN@insertion[label=above:$#3$]{tnX}{(0.80,0)}
        \TN@tensor{tnR}{(1.60,0)}
        \TN@pleg{tnL}{north}{0,\TN@physleglen}
        \TN@pleg{tnR}{north}{0,\TN@physleglen}
        \TN@uplabel{tnL}{\TN@physlabeloff}{#1}
        \TN@uplabel{tnR}{\TN@physlabeloff}{#2}
        \TN@vconnect{tnL}{east}{tnX}{west}
        \TN@vconnect{tnX}{east}{tnR}{west}
        \TN@vtracebelow{tnL}{tnR}{0.78}
        \node at (0.80,-0.88) {$\operatorname{Tr}$};
    \end{tikzpicture}%
}

\newcommand{\TNExternalTraceInsertion}[3]{%
    \begin{tikzpicture}[tn picture]
        \TN@tensor{tnL}{(0,0)}
        \TN@tensor{tnR}{(1.60,0)}
        \TN@insertion[label=above:$#3$]{tnY}{(0.80,-0.78)}
        \TN@pleg{tnL}{north}{0,\TN@physleglen}
        \TN@pleg{tnR}{north}{0,\TN@physleglen}
        \TN@uplabel{tnL}{\TN@physlabeloff}{#1}
        \TN@uplabel{tnR}{\TN@physlabeloff}{#2}
        \TN@vconnect{tnL}{east}{tnR}{west}
        \TN@vtracevia{tnL}{tnY}{tnR}{0.15}{0.15}
        \node at (0.80,-1.00) {$\operatorname{Tr}$};
    \end{tikzpicture}%
}

\newcommand{\TNBoundaryRegrow}[4]{%
    \begin{tikzpicture}[tn picture]
        \TN@tensor{rA}{(0,0)}
        \TN@tensor{rB}{(1.05,0)}
        \TN@tensor{rC}{(2.10,0)}
        \TN@insertion[label=above:$#1$]{rX}{(3.55,0)}
        \foreach \n in {rA,rB,rC} {
                \TN@pleg{\n}{north}{0,0.44}
            }
        \TN@uplabel{rA}{0.66}{#2}
        \TN@uplabel{rC}{0.66}{#3}
        \TN@vleg{rA}{west}{-0.50,0}
        \TN@vconnect{rA}{east}{rB}{west}
        \TN@vconnect{rB}{east}{rC}{west}
        \TN@vconnect{rC}{east}{rX}{west}
        \TN@vleg{rX}{east}{0.45,0}
        \node at (1.78,-0.58) {$#4$};
    \end{tikzpicture}%
}

\newcommand{\TNLocalEqualityStep}[3]{%
    \begin{tikzpicture}[tn picture]
        \TN@tensor{lA}{(0,0)}
        \TN@insertion[label=above:$#1$]{lX}{(0.82,0)}
        \TN@tensor{lB}{(1.65,0)}
        \TN@pleg{lB}{north}{0,0.44}
        \TN@uplabel{lB}{0.66}{#3}
        \TN@vleg{lA}{west}{-0.50,0}
        \TN@vconnect{lA}{east}{lX}{west}
        \TN@vconnect{lX}{east}{lB}{west}
        \TN@vleg{lB}{east}{0.50,0}
        \node at (2.85,0) {$=$};
        \TN@tensor{rA}{(4.10,0)}
        \TN@insertion[label=above:$#2$]{rX}{(4.92,0)}
        \TN@tensor{rB}{(5.75,0)}
        \TN@pleg{rB}{north}{0,0.44}
        \TN@uplabel{rB}{0.66}{#3}
        \TN@vleg{rA}{west}{-0.50,0}
        \TN@vconnect{rA}{east}{rX}{west}
        \TN@vconnect{rX}{east}{rB}{west}
        \TN@vleg{rB}{east}{0.50,0}
    \end{tikzpicture}%
}

\newcommand{\TNGroundSpaceMap}[5]{%
    \begin{tikzpicture}[tn picture]
        \TN@tensor{tnL}{(0,0)}
        \TN@tensor{tnM}{(1.15,0)}
        \TN@tensor{tnR}{(2.95,0)}
        \TN@pleg{tnL}{north}{0,\TN@physleglen}
        \TN@pleg{tnM}{north}{0,\TN@physleglen}
        \TN@pleg{tnR}{north}{0,\TN@physleglen}
        \TN@uplabel{tnL}{\TN@physlabeloff}{#2}
        \TN@uplabel{tnR}{\TN@physlabeloff}{#3}
        \TN@vleg{tnL}{west}{-0.45,0}
        \TN@vconnect{tnL}{east}{tnM}{west}
        \TN@vleg{tnM}{east}{0.56,0}
        \TN@vleg{tnR}{west}{-0.56,0}
        \TN@vleg{tnR}{east}{0.45,0}
        \node at (2.05,0) {$\cdots$};
        \TN@insertion[label=below:$#5$]{tnX}{(1.48,-0.82)}
        \TN@vtracevia{tnL}{tnX}{tnR}{0.45}{0.45}
        \node at ($(tnM.south)+(0,-0.28)$) {$#1$};
        \node at (2.05,-0.55) {$#4$};
    \end{tikzpicture}%
}

% ---------- PEPS diagrams ----------

\newcommand{\TNPEPSEdgeBlockingReduction}{%
    \begin{tikzpicture}[tn picture]
        \begin{scope}
            \TN@pepsFiveSitePatch{blockC}
            \TN@vpath[tn selected path]{(blockCone) circle (0.30cm)}
            \TN@vpath[tn selected path]{(blockCfive) circle (0.30cm)}
            \TN@vpath[tn selected path, rounded corners=2mm]
                {(0.40,-0.20) rectangle (1.70,1.70)}
            \node[tn region label selected] at (-0.45,-0.35) {$A'_2$};
            \node[tn region label selected] at (-1.10,1.80) {$A'_1$};
            \node[tn region label selected] at (1.70,1.80) {$A'_3$};
        \end{scope}
        \node at (2.70,0.70) {$\rightsquigarrow$};
        \begin{scope}[xshift=3.75cm,yshift=0.70cm]
            \TN@factorpath{(0.20,0) rectangle (4.00,-0.46)}
            \TN@tensor{blkL}{(0.80,0)}
            \TN@tensor{blkM}{(2.10,0)}
            \TN@tensor{blkR}{(3.40,0)}
            \foreach \n in {blkL,blkM,blkR} {
                    \TN@pleg{\n}{north}{0,0.46}
                }
            \node at ($(blkL.south)+(0,-0.30)$) {$A'_1$};
            \node at ($(blkM.south)+(0,-0.30)$) {$A'_3$};
            \node at ($(blkR.south)+(0,-0.30)$) {$A'_2$};
        \end{scope}
    \end{tikzpicture}%
}

\newcommand{\TNPEPSEdgeInsertedCoeff}{%
    \begin{tikzpicture}[tn picture]
        \TN@boxedThreeSiteInsertion{edgeCoeff}{A'_1}{A'_3}{A'_2}{X}
    \end{tikzpicture}%
}

\newcommand{\TNPEPSThreeSiteInsertionComparison}{%
    \begin{tikzpicture}[tn picture]
        \begin{scope}
            \TN@boxedThreeSiteInsertion{insA}{A_1}{A_2}{A_3}{X}
        \end{scope}
        \node at (4.20,0) {$=$};
        \begin{scope}[xshift=4.90cm]
            \TN@boxedThreeSiteInsertion{insB}{B_1}{B_2}{B_3}{Y}
        \end{scope}
    \end{tikzpicture}%
}

\newcommand{\TNPEPSInsertionPhysicalRealization}{%
    \begin{tikzpicture}[tn picture]
        \begin{scope}
            \TN@tensor{realLeftA}{(0,0)}
            \TN@pleg{realLeftA}{north}{0,0.50}
            \TN@vleg{realLeftA}{east}{0.70,0}
            \TN@vleg{realLeftA}{south}{0,-0.58}
            \TN@insertion[label=above:$M$]{realLeftM}{(-0.56,0)}
            \TN@vleg{realLeftM}{west}{-0.42,0}
            \TN@vconnect{realLeftM}{east}{realLeftA}{west}
            \node at ($(realLeftA.south)+(0,-0.30)$) {$A_u$};
            \node at (-0.56,-0.35) {$e$};
            \node at (0.72,-0.36) {$\eta_{\widehat e}$};
            \node at ($(realLeftA.north)+(0,0.73)$) {$i$};
        \end{scope}
        \node at (1.80,0) {$=$};
        \begin{scope}[xshift=3.20cm]
            \TN@tensor{realRightA}{(0,0)}
            \TN@pleg{realRightA}{north}{0,0.30}
            \TN@insertion[label=left:$O$]{realRightO}{(0,0.62)}
            \TN@pconnect{realRightA}{north}{realRightO}{south}
            \TN@pleg{realRightO}{north}{0,0.28}
            \TN@vleg{realRightA}{west}{-0.70,0}
            \TN@vleg{realRightA}{east}{0.70,0}
            \TN@vleg{realRightA}{south}{0,-0.58}
            \node at ($(realRightA.south)+(0,-0.30)$) {$A_u$};
            \node at (-0.72,-0.36) {$\eta_e$};
            \node at (0.72,-0.36) {$\eta_{\widehat e}$};
            \node at ($(realRightO.north)+(0,0.55)$) {$O(A_u(\eta_e,\eta_{\widehat e}))$};
        \end{scope}
    \end{tikzpicture}%
}

\newcommand{\TNPEPSPhysicalToVirtualInsertion}{%
    \begin{tikzpicture}[tn picture]
        \begin{scope}
            \TN@boxedThreeSitePhysicalOne{physVirtL}{B_1}{B_2}{B_3}{O_1}
        \end{scope}
        \node at (3.75,0) {$=$};
        \begin{scope}[xshift=4.15cm]
            \TN@boxedThreeSitePhysicalTwo{physVirtR}{B_1}{B_2}{B_3}{O_2}
        \end{scope}
        \node at (7.95,0) {$\Longrightarrow$};
        \begin{scope}[xshift=8.40cm]
            \TN@boxedThreeSiteInsertion{physVirtW}{B_1}{B_2}{B_3}{W}
        \end{scope}
    \end{tikzpicture}%
}

\newcommand{\TNPEPSEdgeInsertionEquality}{%
    \begin{tikzpicture}[tn picture]
        \begin{scope}
            \TN@pepsFiveSitePatch{edgeEqA}
            \TN@pepsFiveSitePatchLabels{edgeEqA}{A_1}{A_2}{A_3}{A_4}{A_5}
            \TN@insertion[label=right:$X$]{edgeInsA}
                {($(edgeEqAtwo)!0.50!(edgeEqAfive)$)}
        \end{scope}
        \node at (2.75,0.70) {$=$};
        \begin{scope}[xshift=4.15cm]
            \TN@pepsFiveSitePatch{edgeEqB}
            \TN@pepsFiveSitePatchLabels{edgeEqB}
            {\widetilde B_1}{\widetilde B_2}{\widetilde B_3}
            {\widetilde B_4}{\widetilde B_5}
            \TN@insertion[label=right:$X$]{edgeInsB}
                {($(edgeEqBtwo)!0.50!(edgeEqBfive)$)}
        \end{scope}
    \end{tikzpicture}%
}

\newcommand{\TNPEPSEdgeGaugeAbsorption}{%
    \begin{tikzpicture}[tn picture]
        \begin{scope}
            \TN@tensor{absB}{(0,0)}
            \TN@pleg{absB}{north}{0,0.44}
            \TN@vleg{absB}{west}{-0.62,0}
            \TN@vleg{absB}{east}{0.62,0}
            \TN@vleg{absB}{225}{-0.50,-0.50}
            \TN@vleg{absB}{315}{0.50,-0.50}
            \node at ($(absB.south)+(0,-0.95)$) {$B_v$};
        \end{scope}
        \node at (1.15,0) {$\longmapsto$};
        \begin{scope}[xshift=2.30cm]
            \TN@tensor{absBt}{(0,0)}
            \TN@pleg{absBt}{north}{0,0.44}
            \TN@insertion[label=above:$Z_{e_1}^{\pm1}$]{absL}{(-0.70,0)}
            \TN@insertion[label=above:$Z_{e_2}^{\pm1}$]{absR}{(0.70,0)}
            \TN@insertion[label=left:$Z_{e_3}^{\pm1}$]{absDL}{(-0.48,-0.60)}
            \TN@insertion[label=right:$Z_{e_4}^{\pm1}$]{absDR}{(0.48,-0.60)}
            \TN@vleg{absL}{west}{-0.28,0}
            \TN@vconnect{absL}{east}{absBt}{west}
            \TN@vconnect{absBt}{east}{absR}{west}
            \TN@vleg{absR}{east}{0.28,0}
            \TN@vconnect{absBt}{225}{absDL}{north}
            \TN@vconnect{absBt}{315}{absDR}{north}
            \TN@vleg{absDL}{south}{-0.20,-0.28}
            \TN@vleg{absDR}{south}{0.20,-0.28}
            \node at ($(absBt.south)+(0,-1.10)$) {$\widetilde B_v$};
        \end{scope}
    \end{tikzpicture}%
}

\newcommand{\TNPEPSTwoInjectiveTensorInsertionComparison}{%
    \begin{tikzpicture}[tn picture]
        \begin{scope}
            \TN@twoInjectiveInsertionPair{twoInjA}{A_1}{A_2}{X}
        \end{scope}
        \node at (3.20,0) {$=$};
        \begin{scope}[xshift=3.90cm]
            \TN@twoInjectiveInsertionPair{twoInjB}{B_1}{B_2}{X}
        \end{scope}
    \end{tikzpicture}%
}

\newcommand{\TNPEPSTwoInjectiveGaugeScalarReduction}{%
    \begin{tikzpicture}[tn picture]
        \begin{scope}
            \TN@threeLegTensorFan{gaugeRedA}{A_1}
        \end{scope}
        \node at (1.25,0) {$=$};
        \begin{scope}[xshift=2.05cm]
            \TN@threeLegTensorFan{gaugeRedBZ}{B_1}
            \TN@insertion[label=right:$Z$]{gaugeRedZ}{(0.52,-0.36)}
        \end{scope}
        \node at (3.30,0) {$=$};
        \begin{scope}[xshift=4.10cm]
            \TN@threeLegTensorFan{gaugeRedBU}{B_1}
            \TN@insertion[label=right:$U$]{gaugeRedU}{(0.56,0.30)}
        \end{scope}
        \node at (5.35,0) {$=$};
        \begin{scope}[xshift=6.15cm]
            \TN@threeLegTensorFan{gaugeRedBW}{B_1}
            \TN@insertion[label=right:$W$]{gaugeRedW}{(0.52,-0.04)}
        \end{scope}
    \end{tikzpicture}%
}

\newcommand{\TNPEPSOneVertexComplementComparison}{%
    \begin{tikzpicture}[tn picture]
        \begin{scope}
            \TN@tensor{oneSiteAR}{(1.00,0)}
            \TN@tensor{oneSiteA}{(2.00,0)}
            \TN@pleg{oneSiteAR}{north}{0,0.42}
            \TN@pleg{oneSiteA}{north}{0,0.42}
            \TN@vleg{oneSiteAR}{west}{-0.50,0}
            \TN@vconnect{oneSiteAR}{east}{oneSiteA}{west}
            \TN@vleg{oneSiteA}{east}{0.50,0}
            \node at ($(oneSiteAR.south)+(0,-0.30)$) {$A_R$};
            \node at ($(oneSiteA.south)+(0,-0.30)$) {$A_v$};
        \end{scope}
        \node at (3.12,0) {$=$};
        \begin{scope}[xshift=3.55cm]
            \TN@tensor{oneSiteAS}{(1.00,0)}
            \TN@pleg{oneSiteAS}{north}{0,0.42}
            \TN@vleg{oneSiteAS}{west}{-0.50,0}
            \TN@vleg{oneSiteAS}{east}{0.50,0}
            \node at ($(oneSiteAS.south)+(0,-0.30)$) {$A_S$};
        \end{scope}
        \node at (5.42,0) {$\propto$};
        \begin{scope}[xshift=5.95cm]
            \TN@tensor{oneSiteBS}{(1.00,0)}
            \TN@pleg{oneSiteBS}{north}{0,0.42}
            \TN@vleg{oneSiteBS}{west}{-0.50,0}
            \TN@vleg{oneSiteBS}{east}{0.50,0}
            \node at ($(oneSiteBS.south)+(0,-0.30)$) {$\widetilde B_S$};
        \end{scope}
        \node at (7.82,0) {$=$};
        \begin{scope}[xshift=8.25cm]
            \TN@tensor{oneSiteBR}{(1.00,0)}
            \TN@tensor{oneSiteB}{(2.00,0)}
            \TN@pleg{oneSiteBR}{north}{0,0.42}
            \TN@pleg{oneSiteB}{north}{0,0.42}
            \TN@vleg{oneSiteBR}{west}{-0.50,0}
            \TN@vconnect{oneSiteBR}{east}{oneSiteB}{west}
            \TN@vleg{oneSiteB}{east}{0.50,0}
            \node at ($(oneSiteBR.south)+(0,-0.30)$) {$\widetilde B_R$};
            \node at ($(oneSiteB.south)+(0,-0.30)$) {$\widetilde B_v$};
        \end{scope}
    \end{tikzpicture}%
}

\newcommand{\TNPEPSInjectiveRegionUnion}{%
    \begin{tikzpicture}[tn picture]
        \TN@tensor{unionAB}{(0,0)}
        \TN@tensor{unionI}{(0.70,0.70)}
        \TN@tensor{unionBA}{(2.00,0)}
        \TN@tensor{unionC}{(1.30,-0.70)}
        \foreach \n in {unionAB,unionI,unionBA,unionC} {
                \TN@pleg{\n}{north}{0,0.42}
            }
        \TN@vport{unionABToI}{unionAB}{45}
        \TN@vport{unionIToAB}{unionI}{225}
        \TN@vport{unionIToBA}{unionI}{-28}
        \TN@vport{unionBAToI}{unionBA}{152}
        \TN@vport{unionBAToC}{unionBA}{225}
        \TN@vport{unionCToBA}{unionC}{45}
        \TN@vport{unionCToAB}{unionC}{152}
        \TN@vport{unionABToC}{unionAB}{-28}
        \TN@vport{unionABToBA}{unionAB}{east}
        \TN@vport{unionBAToAB}{unionBA}{west}
        \TN@vport{unionIToC}{unionI}{-67}
        \TN@vport{unionCToI}{unionC}{113}
        \TN@vconnectports{unionABToI}{unionIToAB}
        \TN@vconnectports{unionIToBA}{unionBAToI}
        \TN@vconnectports{unionBAToC}{unionCToBA}
        \TN@vconnectports{unionCToAB}{unionABToC}
        \TN@vconnectports{unionABToBA}{unionBAToAB}
        \TN@vconnectports{unionIToC}{unionCToI}
        \node[anchor=east] at ($(unionAB)+(-0.12,-0.20)$) {$A\setminus B$};
        \node[anchor=west] at ($(unionI)+(0.12,0.02)$) {$A\cap B$};
        \node[anchor=west] at ($(unionBA)+(0.12,-0.20)$) {$B\setminus A$};
        \node at ($(unionC)+(0,-0.38)$) {$(A\cup B)^c$};
    \end{tikzpicture}%
}

\newcommand{\TNPEPSInjectiveRegionUnionProof}{%
    \begin{tikzpicture}[tn picture, scale=0.82]
        \begin{scope}
            \TN@tensor{proofAB}{(0,0)}
            \TN@tensor{proofI}{(0.70,0.70)}
            \TN@tensor{proofBA}{(2.00,0)}
            \TN@insertion{proofX}{(1.30,-0.70)}
            \foreach \n in {proofAB,proofI,proofBA} {
                    \TN@pleg{\n}{north}{0,0.36}
                }
            \TN@vport{proofABToI}{proofAB}{45}
            \TN@vport{proofIToAB}{proofI}{225}
            \TN@vport{proofIToBA}{proofI}{-28}
            \TN@vport{proofBAToI}{proofBA}{152}
            \TN@vport{proofBAToX}{proofBA}{225}
            \TN@vport{proofXToBA}{proofX}{45}
            \TN@vport{proofXToAB}{proofX}{152}
            \TN@vport{proofABToX}{proofAB}{-28}
            \TN@vport{proofABToBA}{proofAB}{east}
            \TN@vport{proofBAToAB}{proofBA}{west}
            \TN@vport{proofIToX}{proofI}{-67}
            \TN@vport{proofXToI}{proofX}{113}
            \TN@vconnectports{proofABToI}{proofIToAB}
            \TN@vconnectports{proofIToBA}{proofBAToI}
            \TN@vconnectports{proofBAToX}{proofXToBA}
            \TN@vconnectports{proofXToAB}{proofABToX}
            \TN@vconnectports{proofABToBA}{proofBAToAB}
            \TN@vconnectports{proofIToX}{proofXToI}
            \node at ($(proofX)+(0,-0.34)$) {$X$};
            \node at (1.00,-1.18) {$=0$};
        \end{scope}
        \TN@arrow{(2.60,0.00) -- (3.20,0.00)
            node[midway, above] {$A^{-1}$}}
        \begin{scope}[xshift=3.80cm]
            \coordinate (proofAone) at (0,0);
            \coordinate (proofItwo) at (0.70,0.70);
            \TN@tensor{proofBAone}{(2.00,0)}
            \TN@insertion{proofXone}{(1.30,-0.70)}
            \TN@pleg{proofBAone}{north}{0,0.36}
            \TN@vport{proofBAoneToX}{proofBAone}{225}
            \TN@vport{proofXoneToBA}{proofXone}{45}
            \TN@vport{proofXoneToA}{proofXone}{152}
            \TN@vport{proofXoneToI}{proofXone}{113}
            \TN@vport{proofBAoneToA}{proofBAone}{west}
            \TN@vport{proofBAoneToI}{proofBAone}{152}
            \TN@vterminal{proofXoneToAEnd}{($(proofXone)!0.34!(proofAone)$)}
            \TN@vterminal{proofXoneToIEnd}{($(proofXone)!0.34!(proofItwo)$)}
            \TN@vterminal{proofBAoneToAEnd}{($(proofBAone)!0.34!(proofAone)$)}
            \TN@vterminal{proofBAoneToIEnd}{($(proofBAone)!0.34!(proofItwo)$)}
            \TN@vconnectports{proofBAoneToX}{proofXoneToBA}
            \TN@vconnectports{proofXoneToA}{proofXoneToAEnd}
            \TN@vconnectports{proofXoneToI}{proofXoneToIEnd}
            \TN@vconnectports{proofBAoneToA}{proofBAoneToAEnd}
            \TN@vconnectports{proofBAoneToI}{proofBAoneToIEnd}
            \node at ($(proofXone)+(0,-0.34)$) {$X$};
            \node at (1.00,-1.18) {$=0$};
        \end{scope}
        \TN@arrow{(6.45,0.00) -- (7.05,0.00)
            node[midway, above] {$A\cap B$}}
        \begin{scope}[xshift=7.70cm]
            \coordinate (proofAtwo) at (0,0);
            \TN@tensor{proofItwoNode}{(0.70,0.70)}
            \TN@tensor{proofBAtwo}{(2.00,0)}
            \TN@insertion{proofXtwo}{(1.30,-0.70)}
            \foreach \n in {proofItwoNode,proofBAtwo} {
                    \TN@pleg{\n}{north}{0,0.36}
                }
            \TN@vport{proofBAtwoToX}{proofBAtwo}{225}
            \TN@vport{proofXtwoToBA}{proofXtwo}{45}
            \TN@vport{proofXtwoToI}{proofXtwo}{113}
            \TN@vport{proofItwoNodeToX}{proofItwoNode}{-67}
            \TN@vport{proofItwoNodeToBA}{proofItwoNode}{-28}
            \TN@vport{proofBAtwoToI}{proofBAtwo}{152}
            \TN@vport{proofXtwoToA}{proofXtwo}{152}
            \TN@vport{proofBAtwoToA}{proofBAtwo}{west}
            \TN@vport{proofItwoNodeToA}{proofItwoNode}{225}
            \TN@vterminal{proofXtwoToAEnd}{($(proofXtwo)!0.34!(proofAtwo)$)}
            \TN@vterminal{proofBAtwoToAEnd}{($(proofBAtwo)!0.34!(proofAtwo)$)}
            \TN@vterminal{proofItwoToAEnd}{
                ($(proofItwoNode)!0.34!(proofAtwo)$)}
            \TN@vconnectports{proofBAtwoToX}{proofXtwoToBA}
            \TN@vconnectports{proofXtwoToI}{proofItwoNodeToX}
            \TN@vconnectports{proofItwoNodeToBA}{proofBAtwoToI}
            \TN@vconnectports{proofXtwoToA}{proofXtwoToAEnd}
            \TN@vconnectports{proofBAtwoToA}{proofBAtwoToAEnd}
            \TN@vconnectports{proofItwoNodeToA}{proofItwoToAEnd}
            \node at ($(proofXtwo)+(0,-0.34)$) {$X$};
            \node at (1.00,-1.18) {$=0$};
        \end{scope}
        \TN@arrow{(10.35,0.00) -- (10.95,0.00)
            node[midway, above] {$B^{-1}$}}
        \begin{scope}[xshift=11.55cm]
            \coordinate (proofAthree) at (0,0);
            \coordinate (proofIthree) at (0.70,0.70);
            \coordinate (proofBAthree) at (2.00,0);
            \TN@insertion{proofXthree}{(1.30,-0.70)}
            \TN@vport{proofXthreeToA}{proofXthree}{152}
            \TN@vport{proofXthreeToI}{proofXthree}{113}
            \TN@vport{proofXthreeToBA}{proofXthree}{45}
            \TN@vterminal{proofXthreeToAEnd}{($(proofXthree)!0.34!(proofAthree)$)}
            \TN@vterminal{proofXthreeToIEnd}{($(proofXthree)!0.34!(proofIthree)$)}
            \TN@vterminal{proofXthreeToBAEnd}{
                ($(proofXthree)!0.34!(proofBAthree)$)}
            \TN@vconnectports{proofXthreeToA}{proofXthreeToAEnd}
            \TN@vconnectports{proofXthreeToI}{proofXthreeToIEnd}
            \TN@vconnectports{proofXthreeToBA}{proofXthreeToBAEnd}
            \node at ($(proofXthree)+(0,-0.34)$) {$X=0$};
        \end{scope}
    \end{tikzpicture}%
}

\newcommand{\TNPEPSNormalRegionsRS}{%
    \begin{tikzpicture}[tn picture, scale=0.90]
        \begin{scope}
            \clip (-0.25,0.75) rectangle (2.25,3.25);
            \TN@pepsNormalSquareLattice
            \TN@pepsNormalRegionR
        \end{scope}
        \node at (2.85,2.00) {and};
        \begin{scope}[xshift=3.55cm]
            \clip (-0.25,0.75) rectangle (2.25,3.25);
            \TN@pepsNormalSquareLattice
            \TN@pepsNormalRegionS
        \end{scope}
    \end{tikzpicture}%
}

\newcommand{\TNPEPSNormalRegionT}{%
    \begin{tikzpicture}[tn picture, scale=0.90]
        \clip (-0.25,0.25) rectangle (3.25,3.25);
        \TN@pepsNormalSquareLattice
        \TN@pepsNormalRegionT
    \end{tikzpicture}%
}

\newcommand{\TNPEPSNormalRectangleCover}{%
    \begin{tikzpicture}[tn picture, scale=0.78]
        \TN@squarelatticepatch{coverGrid}{4}{4}{0.5}
        \path[tn region selected] (0.3,0.3) rectangle (1.2,1.7);
        \path[tn region secondary] (0.3,0.8) rectangle (1.7,1.7);
        \path[tn region complement, fill=none]
            (0.3,0.3) -- (1.2,0.3) -- (1.2,0.8) -- (1.7,0.8) --
            (1.7,1.7) -- (0.3,1.7) -- cycle;
        \node[tn region label selected] at (0.75,0.55) {$2\times3$};
        \node[tn region label secondary] at (1.20,1.25) {$3\times2$};
        \node[align=center] at (0.98,-0.68)
        {$U=\bigcup_j C_j,$\\[-1pt]
        $C_j\cong 2\times3\ \mathrm{or}\ 3\times2$};
    \end{tikzpicture}%
}

\newcommand{\TNPEPSNormalEdgeComplementTopCollar}{%
    \begin{tikzpicture}[tn picture, scale=0.90]
        \begin{scope}
            \clip (-0.25,-0.25) rectangle (3.25,3.25);
            \TN@pepsNormalSquareLattice
            \TN@pepsNormalRegionT
            \TN@pepsNormalCollar
        \end{scope}
        \node at (3.72,1.48) {$=$};
        \begin{scope}[xshift=4.35cm]
            \clip (-0.25,-0.25) rectangle (3.25,3.25);
            \TN@pepsNormalComplementBlock
            \TN@pepsNormalSquareLattice
        \end{scope}
    \end{tikzpicture}%
}

\newcommand{\TNPEPSNormalOneSiteSeparation}{%
    \begin{tikzpicture}[tn picture, scale=0.90]
        \begin{scope}
            \clip (-0.25,0.75) rectangle (2.25,3.25);
            \TN@pepsNormalSquareLattice
            \TN@pepsNormalRegionR
            \TN@pepsNormalSingleSiteDifference
        \end{scope}
        \node at (2.85,2.00) {$\subset$};
        \begin{scope}[xshift=3.55cm]
            \clip (-0.25,0.75) rectangle (2.25,3.25);
            \TN@pepsNormalSquareLattice
            \TN@pepsNormalRegionS
            \TN@pepsNormalSingleSiteDifference
        \end{scope}
        \node at (4.55,0.52) {$S=R\cup\{v\}$};
    \end{tikzpicture}%
}

\newcommand{\TNPEPSNormalEdgeBlockingReduction}{%
    \begin{tikzpicture}[tn picture, scale=0.90]
        \begin{scope}
            \clip (-0.25,-0.25) rectangle (3.25,3.25);
            \TN@pepsNormalComplementBlock
            \TN@pepsNormalSquareLattice
            \TN@pepsNormalDistinguishedEdge
            \TN@pepsNormalSelectedBlock
            \TN@pepsNormalSecondaryBlock
        \end{scope}
        \node at (3.75,1.45) {$\Rightarrow$};
        \begin{scope}[xshift=4.45cm,yshift=1.50cm]
            \TN@pepsBlockedNormalChain
        \end{scope}
    \end{tikzpicture}%
}

\newcommand{\TNPEPSNormalEdgeBlockingHypotheses}{%
    \begin{tikzpicture}[tn picture, scale=0.90]
        \clip (-0.25,-0.25) rectangle (3.25,3.25);
        \TN@pepsNormalComplementBlock
        \TN@pepsNormalSquareLattice
        \TN@pepsNormalDistinguishedEdge
        \TN@pepsNormalSelectedBlock
        \TN@pepsNormalSecondaryBlock
    \end{tikzpicture}%
}

\newcommand{\TNPEPSNormalBlockingHypotheses}{%
    \begin{tikzpicture}[tn picture, scale=0.72]
        \begin{scope}
            \clip (-0.25,-0.25) rectangle (3.25,3.25);
            \TN@pepsNormalComplementBlock
            \TN@pepsNormalSquareLattice
            \TN@pepsNormalDistinguishedEdge
            \TN@pepsNormalSelectedBlock
            \TN@pepsNormalSecondaryBlock
        \end{scope}
        \node at (1.50,-0.58) {edge blocking};
        \node at (3.75,1.45) {and};
        \begin{scope}[xshift=4.60cm,yshift=-0.20cm]
            \begin{scope}
                \clip (-0.25,0.75) rectangle (2.25,3.25);
                \TN@pepsNormalSquareLattice
                \TN@pepsNormalRegionR
                \TN@pepsNormalSingleSiteDifference
            \end{scope}
            \node at (2.85,2.00) {$\subset$};
            \begin{scope}[xshift=3.55cm]
                \clip (-0.25,0.75) rectangle (2.25,3.25);
                \TN@pepsNormalSquareLattice
                \TN@pepsNormalRegionS
                \TN@pepsNormalSingleSiteDifference
            \end{scope}
            \node at (4.55,0.52) {$S=R\cup\{v\}$};
            \node at (2.85,0.20) {one-site comparison};
        \end{scope}
    \end{tikzpicture}%
}

\newcommand{\TNPEPSTINormalGaugeAbsorption}{%
    \begin{tikzpicture}[tn picture]
        \begin{scope}
            \TN@pepsSquareTensor{tiB}{(0,0)}
            \node at ($(tiB.south)+(0,-0.88)$) {$B$};
        \end{scope}
        \node at (1.25,0) {$=$};
        \begin{scope}[xshift=2.35cm]
            \TN@tensor{tiA}{(0,0)}
            \TN@pleg{tiA}{45}{0.34,0.34}
            \TN@insertion[label=above:$X$]{tiXL}{(-0.76,0)}
            \TN@insertion[label=above:$X^{-1}$]{tiXR}{(0.76,0)}
            \TN@insertion[label=left:$Y$]{tiYD}{(0,-0.76)}
            \TN@insertion[label=left:$Y^{-1}$]{tiYU}{(0,0.76)}
            \TN@vleg{tiXL}{west}{-0.24,0}
            \TN@vconnect{tiXL}{east}{tiA}{west}
            \TN@vconnect{tiA}{east}{tiXR}{west}
            \TN@vleg{tiXR}{east}{0.24,0}
            \TN@vleg{tiYD}{south}{0,-0.24}
            \TN@vconnect{tiYD}{north}{tiA}{south}
            \TN@vconnect{tiA}{north}{tiYU}{south}
            \TN@vleg{tiYU}{north}{0,0.24}
            \node at ($(tiA.south)+(0,-1.20)$) {$\lambda A$};
        \end{scope}
    \end{tikzpicture}%
}

\newcommand{\TNPEPSEdgeGaugeOrientation}{%
    \begin{tikzpicture}[tn picture]
        \TN@pepsvertex{tnU}{(0,0)}{}
        \TN@pepsvertex{tnW}{(3.30,0)}{}
        \TN@insertion[label=above:$X_e$]{tnXl}{(1.00,0)}
        \TN@insertion[label=above:$(X_e^{-1})^\top$]{tnXr}{(2.30,0)}
        \TN@vconnect{tnU}{east}{tnXl}{west}
        \TN@vconnect{tnXl}{east}{tnXr}{west}
        \TN@vconnect{tnXr}{east}{tnW}{west}
        \node at (1.65,0.72) {$e = (u,w)$};
        \node at ($(tnU.south)+(0,-0.28)$) {$u$};
        \node at ($(tnW.south)+(0,-0.28)$) {$w$};
    \end{tikzpicture}%
}

\newcommand{\TNPEPSGaugeVertexAction}{%
    \begin{tikzpicture}[tn picture]
        \TN@tensor{tnV}{(0,0)}
        \TN@pleg{tnV}{45}{0.38,0.38}
        \TN@insertion[label=above:$M_{ve_1}$]{tnL}{(-1.05,0)}
        \TN@insertion[label=above:$M_{ve_2}$]{tnR}{(1.05,0)}
        \TN@insertion[label=left:$M_{ve_3}$]{tnDL}{(-0.78,-0.92)}
        \TN@insertion[label=right:$M_{ve_4}$]{tnDR}{(0.78,-0.92)}
        \TN@vconnect{tnL}{east}{tnV}{west}
        \TN@vconnect{tnV}{east}{tnR}{west}
        \TN@vconnect{tnV}{225}{tnDL}{north}
        \TN@vconnect{tnV}{315}{tnDR}{north}
        \TN@vleg{tnL}{west}{-0.42,0}
        \TN@vleg{tnR}{east}{0.42,0}
        \TN@vleg{tnDL}{south}{-0.30,-0.42}
        \TN@vleg{tnDR}{south}{0.30,-0.42}
        \node at ($(tnV.45)+(0.38,0.50)$) {$\sigma$};
        \node at ($(tnV.south)+(0,-1.45)$) {$A_v \mapsto A_v^X$};
    \end{tikzpicture}%
}

\newcommand{\TNPEPSGaugeCancellation}{%
    \begin{tikzpicture}[tn picture]
        \TN@pepsvertex{tnL}{(0,0)}{}
        \TN@insertion{tnX}{(0.92,0)}
        \TN@insertion{tnXt}{(2.18,0)}
        \TN@pepsvertex{tnR}{(3.10,0)}{}
        \TN@vconnect{tnL}{east}{tnX}{west}
        \TN@vconnect{tnX}{east}{tnXt}{west}
        \TN@vconnect{tnXt}{east}{tnR}{west}
        \node at ($(tnX.north)+(0,0.28)$) {$X_e(j,i)$};
        \node at ($(tnXt.north)+(0,0.28)$) {$((X_e^{-1})^\top)(j,k)$};
        \node at (4.20,0) {$=$};
        \TN@pepsvertex{tnL2}{(5.15,0)}{}
        \TN@pepsvertex{tnR2}{(6.65,0)}{}
        \TN@vconnect{tnL2}{east}{tnR2}{west}
        \node at (5.90,0.54) {$\sum_j X_e(j,i)\,((X_e^{-1})^\top)(j,k) = \delta(i,k)$};
    \end{tikzpicture}%
}

\newcommand{\TNPEPSBlockedMiddleLocalGaugeFormula}{%
    \begin{tikzpicture}[tn picture]
        \TN@tensor{bmA}{(0,0)}
        \TN@tensor{bmB}{(2.55,0)}
        \foreach \a in {10,90,160,220,300} {
                \TN@vleg{bmA}{\a}{\a:0.55}
                \TN@vleg{bmB}{\a}{\a:0.55}
            }
        \TN@pleg{bmA}{45}{0.44,0.44}
        \TN@pleg{bmB}{45}{0.44,0.44}
        \node at ($(bmA)+(0,-0.62)$) {$A_v$};
        \node at ($(bmB)+(0,-0.62)$) {$B_v$};
        \node at (1.28,0.25) {$\prod_{e\ni v}X_e$};
        \TN@arrow{(0.62,0) -- (1.93,0)}
    \end{tikzpicture}%
}

\newcommand{\TNPEPSLocalGaugeExtraction}{%
    \begin{tikzpicture}[tn picture]
        \TN@tensor{tnLv}{(0,0)}
        \TN@pleg{tnLv}{45}{0.38,0.38}
        \node at ($(tnLv.45)+(0.38,0.48)$) {$\sigma$};
        \TN@vleg{tnLv}{west}{-0.90,0}
        \TN@vleg{tnLv}{east}{0.90,0}
        \TN@vleg{tnLv}{225}{-0.72,-0.72}
        \TN@vleg{tnLv}{315}{0.72,-0.72}
        \node at (0,-1.42) {$\mathrm{star}(v)$};
        \node at (2.85,0) {$\xrightarrow[\text{left inverse}]{\text{injective}}$};
        \TN@tensor{tnRv}{(5.70,0)}
        \TN@pleg{tnRv}{45}{0.38,0.38}
        \node at ($(tnRv.45)+(0.38,0.48)$) {$\sigma$};
        \TN@insertion{tnRl}{(4.62,0)}
        \TN@insertion{tnRr}{(6.78,0)}
        \TN@insertion{tnRdl}{(4.95,-0.86)}
        \TN@insertion{tnRdr}{(6.45,-0.86)}
        \TN@vleg{tnRl}{west}{-0.42,0}
        \TN@vconnect{tnRl}{east}{tnRv}{west}
        \TN@vconnect{tnRv}{east}{tnRr}{west}
        \TN@vleg{tnRr}{east}{0.42,0}
        \TN@vconnect{tnRv}{225}{tnRdl}{north}
        \TN@vconnect{tnRv}{315}{tnRdr}{north}
        \node at (5.70,-1.42) {$\text{local edge gauges}$};
    \end{tikzpicture}%
}

\newcommand{\TNPEPSGlobalConsistency}{%
    \begin{tikzpicture}[tn picture]
        \TN@pepsvertex{tnU}{(0,0)}{}
        \TN@pepsvertex{tnW}{(3.35,0)}{}
        \TN@vleg{tnU}{west}{-0.70,0}
        \TN@vleg{tnW}{east}{0.70,0}
        \TN@vleg{tnU}{225}{-0.55,-0.55}
        \TN@vleg{tnW}{315}{0.55,-0.55}
        \TN@insertion[label=above:$X_e$]{tnXl}{(1.02,0)}
        \TN@insertion[label=above:$(X_e^{-1})^\top$]{tnXr}{(2.33,0)}
        \TN@vconnect{tnU}{east}{tnXl}{west}
        \TN@vconnect{tnXl}{east}{tnXr}{west}
        \TN@vconnect{tnXr}{east}{tnW}{west}
        \node at (1.67,0.78) {$\text{shared edge } e$};
        \node at ($(tnU.south)+(0,-0.28)$) {$u$};
        \node at ($(tnW.south)+(0,-0.28)$) {$w$};
    \end{tikzpicture}%
}

% ---------- Pedagogical PEPS overview diagrams ----------

% A PEPS on a square lattice. Each black dot is the local tensor at a vertex;
% the thin horizontal and vertical bonds are the virtual edges of the finite
% graph; the short thick diagonal stub at every vertex is the physical leg.
\newcommand{\TNPEPSLatticeState}{%
    \begin{tikzpicture}[tn picture]
        \TN@squarepepspatch{pl}{2}{2}{1}
    \end{tikzpicture}%
}

% A single local tensor at a vertex $v$: the thick diagonal stub is the
% physical index $\sigma$, the four thin legs are the incident virtual edges.
\newcommand{\TNPEPSLocalTensorStar}{%
    \begin{tikzpicture}[tn picture]
        \TN@pepssiteWithOpenLegs{lt}{(0,0)}{}
        \node at (0.66,0.62) {$\sigma$};
        \node at (-0.98,0) {$e_1$};
        \node at (0.98,0) {$e_2$};
        \node at (0,0.86) {$e_3$};
        \node at (0,-0.86) {$e_4$};
        \node at (-0.40,-0.34) {$A_v$};
    \end{tikzpicture}%
}

% Vertex injectivity. Reading the local tensor as the map sending a virtual
% configuration $\eta$ on the incident edges to the physical vector
% $A_v(\eta,\cdot)\in\C^d$, vertex injectivity is the linear independence of
% this family.
\newcommand{\TNPEPSVertexInjectivityMap}{%
    \begin{tikzpicture}[tn picture]
        \TN@tensor{vi}{(0,0)}
        \TN@pleg{vi}{45}{0.46,0.46}
        \node at (0.72,0.72) {$\C^d$};
        \TN@vleg{vi}{west}{-0.62,0}
        \TN@vleg{vi}{east}{0.62,0}
        \TN@vleg{vi}{225}{-0.46,-0.46}
        \TN@vleg{vi}{315}{0.46,-0.46}
        \node at (0,-1.04) {$\eta$};
        \node at (1.95,0.30) {\scriptsize injective};
        \node at (1.95,0) {$\longmapsto$};
        \node[align=center] at (4.45,0)
        {$\{A_v(\eta,\cdot)\}_{\eta}$ linearly\\ independent in $\C^d$};
    \end{tikzpicture}%
}

% The state coefficient as a contraction. On the five-vertex example graph
% (pentagon with one chord) the physical legs stay open while every virtual
% bond is summed over; the coefficient is the resulting scalar.
\newcommand{\TNPEPSStateContraction}{%
    \begin{tikzpicture}[tn picture]
        \TN@pepsFiveSitePatch{sc}
        \TN@pepsFiveSitePatchLabels{sc}{A_1}{A_2}{A_3}{A_4}{A_5}
    \end{tikzpicture}%
}

% Torus geometry. A finite square-lattice patch whose opposite boundary bonds
% are identified, so the columns and rows are read modulo $n$ and $m$.
\newcommand{\TNPEPSTorusGeometry}{%
    \begin{tikzpicture}[tn picture]
        \TN@squarepepspatch{tg}{2}{2}{1}
        \foreach \y in {0,1,2} {
            \TN@vopenport[tn periodic mark]{tg2\y E}{0.5,0}
            \TN@vopenport[tn periodic mark]{tg0\y W}{-0.5,0}
        }
        \foreach \x in {0,1,2} {
            \TN@vopenport[tn periodic mark]{tg\x2N}{0,0.5}
            \TN@vopenport[tn periodic mark]{tg\x0S}{0,-0.5}
        }
        \node at (1,-1.15) {\scriptsize rows and columns identified modulo $n,m$};
    \end{tikzpicture}%
}

% Vertex balance of edge scalars. The endpoint reading $c_{v,e}$ on each
% incident edge; the gauge is balanced at $v$ when the incident readings
% multiply to one.
\newcommand{\TNPEPSVertexScalarBalance}{%
    \begin{tikzpicture}[tn picture]
        \TN@pepssite{vb}{(0,0)}{}
        \TN@vterminal{vbWOpen}{($(vbW)+(-0.58,0)$)}
        \TN@vterminal{vbEOpen}{($(vbE)+(0.58,0)$)}
        \TN@vterminal{vbNOpen}{($(vbN)+(0,0.58)$)}
        \TN@vterminal{vbSOpen}{($(vbS)+(0,-0.58)$)}
        \TN@pterminal{vbPOpen}{($(vbP)+(0.38,0.38)$)}
        \TN@insertion{vbZW}{(-0.5,0)}
        \TN@insertion{vbZE}{(0.5,0)}
        \TN@insertion{vbZN}{(0,0.5)}
        \TN@insertion{vbZS}{(0,-0.5)}
        \TN@vport{vbZWW}{vbZW}{west}
        \TN@vport{vbZWE}{vbZW}{east}
        \TN@vport{vbZEW}{vbZE}{west}
        \TN@vport{vbZEE}{vbZE}{east}
        \TN@vport{vbZNN}{vbZN}{north}
        \TN@vport{vbZNS}{vbZN}{south}
        \TN@vport{vbZSN}{vbZS}{north}
        \TN@vport{vbZSS}{vbZS}{south}
        \TN@vconnectports{vbW}{vbZWE}
        \TN@vconnectports{vbZWW}{vbWOpen}
        \TN@vconnectports{vbE}{vbZEW}
        \TN@vconnectports{vbZEE}{vbEOpen}
        \TN@vconnectports{vbN}{vbZNS}
        \TN@vconnectports{vbZNN}{vbNOpen}
        \TN@vconnectports{vbS}{vbZSN}
        \TN@vconnectports{vbZSS}{vbSOpen}
        \TN@pconnectports{vbP}{vbPOpen}
        \node at (-0.58,0.28) {$c_{v,e_1}$};
        \node at (0.62,-0.28) {$c_{v,e_2}$};
        \node at (-0.42,0.60) {$c_{v,e_3}$};
        \node at (0.42,-0.60) {$c_{v,e_4}$};
        \node at (2.35,0) {$\displaystyle\prod_{e \ni v} c_{v,e} = 1$};
    \end{tikzpicture}%
}

% ---------- Quantum-channel representation diagrams ----------

% Kraus representation T(rho) = sum_i K_i rho K_i^dagger as a double-layer
% contraction: the Kraus index i is summed between the K_i (ket) and
% K_i^dagger (bra) layers; rho enters on the left, T(rho) leaves on the right.
\newcommand{\TNKrausMap}{%
    \begin{tikzpicture}[tn picture]
        \TN@operatorstate{krRho}{(-1.35,0)}{\rho}
        \TN@doublelayer{kr}{(0,0.44)}{(0,-0.44)}
        \TN@operatorstate{krOut}{(1.55,0)}{T(\rho)}
        \TN@vconnectports{krRhoEUp}{krUW}
        \TN@vconnectports{krRhoEDown}{krDW}
        \TN@vconnectports{krUE}{krOutWUp}
        \TN@vconnectports{krDE}{krOutWDown}
        \node at ($(krUS)!0.5!(krDN)+(0.20,0)$) {$i$};
        \node at ($(krU.north)+(0,0.20)$) {$K_i$};
        \node at ($(krD.south)+(0,-0.20)$) {$K_i^\dagger$};
    \end{tikzpicture}%
}

% Stinespring dilation T(rho) = tr_E[V rho V^dagger]: the isometry V and its
% adjoint map the system into a larger space; the environment E is traced out
% (the loop on the right), leaving the channel output T(rho) on the system legs.
\newcommand{\TNStinespring}{%
    \begin{tikzpicture}[tn picture]
        \TN@operatorstate{sRho}{(-1.35,0)}{\rho}
        \TN@systemancillamapup{sV}{(0,0.5)}{V}
        \TN@systemancillamapdown{sVd}{(0,-0.5)}{V^\dagger}
        \TN@operatorstate{sOut}{(1.55,0)}{T(\rho)}
        \TN@vconnectports{sRhoEUp}{sVW}
        \TN@vconnectports{sRhoEDown}{sVdW}
        \TN@vconnectports{sVE}{sOutWUp}
        \TN@vconnectports{sVdE}{sOutWDown}
        \TN@vtraceportsright{sVAncilla}{sVdAncilla}{2.25}{0.30}
        \node at (2.46,0.98) {$E$};
    \end{tikzpicture}%
}

% Choi--Jamiolkowski isomorphism tau = (T tensor id)(|Omega><Omega|): the channel
% T (double layer) has its input legs bent around the maximally entangled pair
% |Omega|, presenting the map as a matrix on the doubled space.
\newcommand{\TNChoiMatrix}{%
    \begin{tikzpicture}[tn picture]
        \TN@tensor{cmU}{(0,0.42)}
        \TN@tensor{cmD}{(0,-0.42)}
        \TN@pconnect{cmU}{south}{cmD}{north}
        \TN@vleg{cmU}{east}{0.62,0}
        \TN@vleg{cmD}{east}{0.62,0}
        \TN@vbezieropen{cmU}{west}{(-1.0,0.42)}{(-1.0,1.05)}
            {(-0.35,1.05)}{0.97,0}
        \TN@vbezieropen{cmD}{west}{(-1.0,-0.42)}{(-1.0,-1.05)}
            {(-0.35,-1.05)}{0.97,0}
        \node[anchor=east] at (-0.78,0) {$\ket{\Omega}$};
        \node at (0.3,0.66) {$T$};
    \end{tikzpicture}%
}

% Trace-pairing transfer form
% T(rho)=sum_{i,j} t_ij tr(sigma_j rho) sigma_i.  The trace loop passes through
% rho and sigma_j; the scalar coefficient t_ij joins the two basis labels, and
% sigma_i supplies the output matrix legs.
\newcommand{\TNTransferMapTracePairing}{%
    \begin{tikzpicture}[tn picture]
        \TN@component{tpRho}{(-1.55,0)}{\rho}
        \TN@component{tpSigJ}{(-0.48,0)}{\sigma_j}
        \TN@tensor{tpT}{(0.82,0)}
        \TN@label[anchor=north]{tpT.south}{(0,-0.10)}{t_{ij}}
        \TN@component{tpSigI}{(2.12,0)}{\sigma_i}
        \TN@vconnectpoints{tpRho.north east}{tpSigJ.north west}
        \TN@vconnectpoints{tpRho.south east}{tpSigJ.south west}
        \TN@vtraceconnectbezier{tpSigJ}{north east}{(-0.12,0.52)}
            {(-1.92,0.52)}{tpRho}{north west}
        \TN@vtraceconnectbezier{tpSigJ}{south east}{(-0.12,-0.52)}
            {(-1.92,-0.52)}{tpRho}{south west}
        \TN@pconnectbezier{tpSigJ}{north}{(-0.32,0.78)}
            {(0.62,0.78)}{tpT}{north}
        \TN@pconnect{tpT}{east}{tpSigI}{west}
        \TN@vleg{tpSigI}{north east}{0.72,0}
        \TN@vleg{tpSigI}{south east}{0.72,0}
        \node at (-1.08,-0.82) {$\tr(\sigma_j\rho)$};
        \node at (3.22,0) {$T(\rho)$};
    \end{tikzpicture}%
}

% Two-point correlator <X_0 Y_n> = tr(Y E_A^n(X rho^R)): on the transfer-map
% (M_D) line, start from the fixed point rho^R, insert X, propagate by n copies
% of the transfer map E_A, insert Y, and close the line by the trace.
\newcommand{\TNTwoPointCorrelator}{%
    \begin{tikzpicture}[tn picture]
        \TN@component[inner sep=1.6pt]{coR}{(0,0)}{\rho^R}
        \TN@insertion{coX}{(1.05,0)}
        \TN@map{coE}{(2.3,0)}{\E_A^n}
        \TN@insertion{coY}{(3.55,0)}
        \TN@vconnect{coR}{east}{coX}{west}
        \TN@vconnect{coX}{east}{coE}{west}
        \TN@vconnect{coE}{east}{coY}{west}
        \node at (1.05,0.28) {$X$};
        \node at (3.55,0.28) {$Y$};
        \TN@vtracebelow{coR}{coY}{0.85}
        \node at (1.8,-1.08) {$\tr$};
    \end{tikzpicture}%
}

\makeatother
